% Options for packages loaded elsewhere
% Options for packages loaded elsewhere
\PassOptionsToPackage{unicode}{hyperref}
\PassOptionsToPackage{hyphens}{url}
\PassOptionsToPackage{dvipsnames,svgnames,x11names}{xcolor}
%
\documentclass[
  11pt,
  letterpaper,
  numbers=noendperiod,
  parskip=half]{scrartcl}
\usepackage{xcolor}
\usepackage[margin=1in,heightrounded]{geometry}
\usepackage{amsmath,amssymb}
\setcounter{secnumdepth}{3}
\usepackage{iftex}
\ifPDFTeX
  \usepackage[T1]{fontenc}
  \usepackage[utf8]{inputenc}
  \usepackage{textcomp} % provide euro and other symbols
\else % if luatex or xetex
  \usepackage{unicode-math} % this also loads fontspec
  \defaultfontfeatures{Scale=MatchLowercase}
  \defaultfontfeatures[\rmfamily]{Ligatures=TeX,Scale=1}
\fi
\usepackage{libertinus}
\ifPDFTeX\else
  % xetex/luatex font selection
  \setmainfont{Libertinus Serif}
  \setsansfont{Libertinus Sans}
  \setmonofont{JetBrains Mono}[Scale=0.88]
  \setmathfont{Libertinus Math}
\fi
% Use upquote if available, for straight quotes in verbatim environments
\IfFileExists{upquote.sty}{\usepackage{upquote}}{}
\IfFileExists{microtype.sty}{% use microtype if available
  \usepackage[]{microtype}
  \UseMicrotypeSet[protrusion]{basicmath} % disable protrusion for tt fonts
}{}
\makeatletter
\@ifundefined{KOMAClassName}{% if non-KOMA class
  \IfFileExists{parskip.sty}{%
    \usepackage{parskip}
  }{% else
    \setlength{\parindent}{0pt}
    \setlength{\parskip}{6pt plus 2pt minus 1pt}}
}{% if KOMA class
  \KOMAoptions{parskip=half}}
\makeatother
% Make \paragraph and \subparagraph free-standing
\makeatletter
\ifx\paragraph\undefined\else
  \let\oldparagraph\paragraph
  \renewcommand{\paragraph}{
    \@ifstar
      \xxxParagraphStar
      \xxxParagraphNoStar
  }
  \newcommand{\xxxParagraphStar}[1]{\oldparagraph*{#1}\mbox{}}
  \newcommand{\xxxParagraphNoStar}[1]{\oldparagraph{#1}\mbox{}}
\fi
\ifx\subparagraph\undefined\else
  \let\oldsubparagraph\subparagraph
  \renewcommand{\subparagraph}{
    \@ifstar
      \xxxSubParagraphStar
      \xxxSubParagraphNoStar
  }
  \newcommand{\xxxSubParagraphStar}[1]{\oldsubparagraph*{#1}\mbox{}}
  \newcommand{\xxxSubParagraphNoStar}[1]{\oldsubparagraph{#1}\mbox{}}
\fi
\makeatother

\usepackage{color}
\usepackage{fancyvrb}
\newcommand{\VerbBar}{|}
\newcommand{\VERB}{\Verb[commandchars=\\\{\}]}
\DefineVerbatimEnvironment{Highlighting}{Verbatim}{commandchars=\\\{\}}
% Add ',fontsize=\small' for more characters per line
\newenvironment{Shaded}{}{}
\newcommand{\AlertTok}[1]{\textcolor[rgb]{1.00,0.33,0.33}{\textbf{#1}}}
\newcommand{\AnnotationTok}[1]{\textcolor[rgb]{0.42,0.45,0.49}{#1}}
\newcommand{\AttributeTok}[1]{\textcolor[rgb]{0.84,0.23,0.29}{#1}}
\newcommand{\BaseNTok}[1]{\textcolor[rgb]{0.00,0.36,0.77}{#1}}
\newcommand{\BuiltInTok}[1]{\textcolor[rgb]{0.84,0.23,0.29}{#1}}
\newcommand{\CharTok}[1]{\textcolor[rgb]{0.01,0.18,0.38}{#1}}
\newcommand{\CommentTok}[1]{\textcolor[rgb]{0.42,0.45,0.49}{#1}}
\newcommand{\CommentVarTok}[1]{\textcolor[rgb]{0.42,0.45,0.49}{#1}}
\newcommand{\ConstantTok}[1]{\textcolor[rgb]{0.00,0.36,0.77}{#1}}
\newcommand{\ControlFlowTok}[1]{\textcolor[rgb]{0.84,0.23,0.29}{#1}}
\newcommand{\DataTypeTok}[1]{\textcolor[rgb]{0.84,0.23,0.29}{#1}}
\newcommand{\DecValTok}[1]{\textcolor[rgb]{0.00,0.36,0.77}{#1}}
\newcommand{\DocumentationTok}[1]{\textcolor[rgb]{0.42,0.45,0.49}{#1}}
\newcommand{\ErrorTok}[1]{\textcolor[rgb]{1.00,0.33,0.33}{\underline{#1}}}
\newcommand{\ExtensionTok}[1]{\textcolor[rgb]{0.84,0.23,0.29}{\textbf{#1}}}
\newcommand{\FloatTok}[1]{\textcolor[rgb]{0.00,0.36,0.77}{#1}}
\newcommand{\FunctionTok}[1]{\textcolor[rgb]{0.44,0.26,0.76}{#1}}
\newcommand{\ImportTok}[1]{\textcolor[rgb]{0.01,0.18,0.38}{#1}}
\newcommand{\InformationTok}[1]{\textcolor[rgb]{0.42,0.45,0.49}{#1}}
\newcommand{\KeywordTok}[1]{\textcolor[rgb]{0.84,0.23,0.29}{#1}}
\newcommand{\NormalTok}[1]{\textcolor[rgb]{0.14,0.16,0.18}{#1}}
\newcommand{\OperatorTok}[1]{\textcolor[rgb]{0.14,0.16,0.18}{#1}}
\newcommand{\OtherTok}[1]{\textcolor[rgb]{0.44,0.26,0.76}{#1}}
\newcommand{\PreprocessorTok}[1]{\textcolor[rgb]{0.84,0.23,0.29}{#1}}
\newcommand{\RegionMarkerTok}[1]{\textcolor[rgb]{0.42,0.45,0.49}{#1}}
\newcommand{\SpecialCharTok}[1]{\textcolor[rgb]{0.00,0.36,0.77}{#1}}
\newcommand{\SpecialStringTok}[1]{\textcolor[rgb]{0.01,0.18,0.38}{#1}}
\newcommand{\StringTok}[1]{\textcolor[rgb]{0.01,0.18,0.38}{#1}}
\newcommand{\VariableTok}[1]{\textcolor[rgb]{0.89,0.38,0.04}{#1}}
\newcommand{\VerbatimStringTok}[1]{\textcolor[rgb]{0.01,0.18,0.38}{#1}}
\newcommand{\WarningTok}[1]{\textcolor[rgb]{1.00,0.33,0.33}{#1}}

\usepackage{longtable,booktabs,array}
\newcounter{none} % for unnumbered tables
\usepackage{calc} % for calculating minipage widths
% Correct order of tables after \paragraph or \subparagraph
\usepackage{etoolbox}
\makeatletter
\patchcmd\longtable{\par}{\if@noskipsec\mbox{}\fi\par}{}{}
\makeatother
% Allow footnotes in longtable head/foot
\IfFileExists{footnotehyper.sty}{\usepackage{footnotehyper}}{\usepackage{footnote}}
\makesavenoteenv{longtable}
\usepackage{graphicx}
\usepackage{tikz}
\usepackage{pgfplots}
\usepgfplotslibrary{groupplots}
\usepackage{tabularx}
\usepackage{makecell}
\usepackage{adjustbox}
\usepackage{siunitx}
\usepackage{colortbl}
\usepackage{enumitem}
\makeatletter
\newsavebox\pandoc@box
\newcommand*\pandocbounded[1]{% scales image to fit in text height/width
  \sbox\pandoc@box{#1}%
  \Gscale@div\@tempa{\textheight}{\dimexpr\ht\pandoc@box+\dp\pandoc@box\relax}%
  \Gscale@div\@tempb{\linewidth}{\wd\pandoc@box}%
  \ifdim\@tempb\p@<\@tempa\p@\let\@tempa\@tempb\fi% select the smaller of both
  \ifdim\@tempa\p@<\p@\scalebox{\@tempa}{\usebox\pandoc@box}%
  \else\usebox{\pandoc@box}%
  \fi%
}
% Set default figure placement to htbp
\def\fps@figure{htbp}
\makeatother


% definitions for citeproc citations
\NewDocumentCommand\citeproctext{}{}
\NewDocumentCommand\citeproc{mm}{%
  \begingroup\def\citeproctext{#2}\cite{#1}\endgroup}
\makeatletter
 % allow citations to break across lines
 \let\@cite@ofmt\@firstofone
 % avoid brackets around text for \cite:
 \def\@biblabel#1{}
 \def\@cite#1#2{{#1\if@tempswa , #2\fi}}
\makeatother
\newlength{\cslhangindent}
\setlength{\cslhangindent}{1.5em}
\newlength{\csllabelwidth}
\setlength{\csllabelwidth}{3em}
\newenvironment{CSLReferences}[2] % #1 hanging-indent, #2 entry-spacing
 {\begin{list}{}{%
  \setlength{\itemindent}{0pt}
  \setlength{\leftmargin}{0pt}
  \setlength{\parsep}{0pt}
  % turn on hanging indent if param 1 is 1
  \ifodd #1
   \setlength{\leftmargin}{\cslhangindent}
   \setlength{\itemindent}{-1\cslhangindent}
  \fi
  % set entry spacing
  \setlength{\itemsep}{#2\baselineskip}}}
 {\end{list}}
\usepackage{calc}
\newcommand{\CSLBlock}[1]{\hfill\break\parbox[t]{\linewidth}{\strut\ignorespaces#1\strut}}
\newcommand{\CSLLeftMargin}[1]{\parbox[t]{\csllabelwidth}{\strut#1\strut}}
\newcommand{\CSLRightInline}[1]{\parbox[t]{\linewidth - \csllabelwidth}{\strut#1\strut}}
\newcommand{\CSLIndent}[1]{\hspace{\cslhangindent}#1}



\setlength{\emergencystretch}{3em} % prevent overfull lines

\providecommand{\tightlist}{%
  \setlength{\itemsep}{0pt}\setlength{\parskip}{0pt}}



 


\usepackage{microtype}
\usepackage{hyperref}
\hypersetup{
  pdftitle={Context Routing for AI Coding Agents},
  pdfauthor={Gunther Brunner},
  pdfsubject={Context allocation policies for AI coding agents},
  pdfkeywords={AI agents, context routing, context engineering, retrieval, vendoring, MCP, agent skills}
}
\usepackage{titling}
\pretitle{\begin{center}\fontsize{22pt}{26pt}\selectfont\bfseries}
\posttitle{\par\end{center}\vskip 0.5em}
\preauthor{\begin{center}\large}
\postauthor{\end{center}}
\predate{\begin{center}\large}
\postdate{\par\end{center}\vskip 1em}
\usepackage[automark,headsepline,footsepline]{scrlayer-scrpage}
\clearpairofpagestyles
\ihead{\small Context Routing for AI Coding Agents}
\ifoot{\small Brunner, 2026}
\ofoot{\pagemark}
\KOMAoption{captions}{tableheading}
\makeatletter
\@ifpackageloaded{tcolorbox}{}{\usepackage[skins,breakable]{tcolorbox}}
\@ifpackageloaded{fontawesome5}{}{\usepackage{fontawesome5}}
\definecolor{paper-callout-color}{HTML}{909090}
\definecolor{paper-callout-note-color}{HTML}{0758E5}
\definecolor{paper-callout-important-color}{HTML}{CC1914}
\definecolor{paper-callout-warning-color}{HTML}{EB9113}
\definecolor{paper-callout-tip-color}{HTML}{00A047}
\definecolor{paper-callout-caution-color}{HTML}{FC5300}
\definecolor{paper-callout-color-frame}{HTML}{acacac}
\definecolor{paper-callout-note-color-frame}{HTML}{4582ec}
\definecolor{paper-callout-important-color-frame}{HTML}{d9534f}
\definecolor{paper-callout-warning-color-frame}{HTML}{f0ad4e}
\definecolor{paper-callout-tip-color-frame}{HTML}{02b875}
\definecolor{paper-callout-caution-color-frame}{HTML}{fd7e14}
\makeatother
\makeatletter
\@ifpackageloaded{caption}{}{\usepackage{caption}}
\AtBeginDocument{%
\ifdefined\contentsname
  \renewcommand*\contentsname{Table of contents}
\else
  \newcommand\contentsname{Table of contents}
\fi
\ifdefined\listfigurename
  \renewcommand*\listfigurename{List of Figures}
\else
  \newcommand\listfigurename{List of Figures}
\fi
\ifdefined\listtablename
  \renewcommand*\listtablename{List of Tables}
\else
  \newcommand\listtablename{List of Tables}
\fi
\ifdefined\figurename
  \renewcommand*\figurename{Figure}
\else
  \newcommand\figurename{Figure}
\fi
\ifdefined\tablename
  \renewcommand*\tablename{Table}
\else
  \newcommand\tablename{Table}
\fi
}
\@ifpackageloaded{float}{}{\usepackage{float}}
\floatstyle{ruled}
\@ifundefined{c@chapter}{\newfloat{codelisting}{h}{lop}}{\newfloat{codelisting}{h}{lop}[chapter]}
\floatname{codelisting}{Listing}
\newcommand*\listoflistings{\listof{codelisting}{List of Listings}}
\makeatother
\makeatletter
\makeatother
\makeatletter
\@ifpackageloaded{caption}{}{\usepackage{caption}}
\@ifpackageloaded{subcaption}{}{\usepackage{subcaption}}
\makeatother
\makeatletter
\@ifpackageloaded{tcolorbox}{}{\usepackage[skins,breakable]{tcolorbox}}
\makeatother
\makeatletter
\@ifundefined{shadecolor}{\definecolor{shadecolor}{HTML}{4a90d9}}{}
\makeatother
\makeatletter
\@ifundefined{codebgcolor}{\definecolor{codebgcolor}{HTML}{f5f5f5}}{}
\makeatother
\makeatletter
\ifdefined\Shaded\renewenvironment{Shaded}{\begin{tcolorbox}[borderline west={3pt}{0pt}{shadecolor}, boxrule=0pt, breakable, colback={codebgcolor}, enhanced, frame hidden, sharp corners]}{\end{tcolorbox}}\fi
\makeatother
\usepackage{bookmark}
\IfFileExists{xurl.sty}{\usepackage{xurl}}{} % add URL line breaks if available
\urlstyle{same}
\hypersetup{
  pdftitle={Context Routing for AI Coding Agents},
  pdfauthor={Gunther Brunner},
  pdfkeywords={AI coding agents, context routing, context
engineering, retrieval-augmented code generation, repository-level code
retrieval, vendoring, Model Context Protocol, agent skills, software
engineering tools},
  colorlinks=true,
  linkcolor={RoyalBlue},
  filecolor={Maroon},
  citecolor={RoyalBlue},
  urlcolor={RoyalBlue},
  pdfcreator={LaTeX}}

\usetikzlibrary{
  arrows.meta,
  backgrounds,
  calc,
  decorations.pathreplacing,
  fit,
  matrix,
  patterns.meta,
  positioning,
  shapes.geometric
}
\pgfplotsset{compat=1.18}

\definecolor{Ink}{HTML}{171717}
\definecolor{MutedInk}{HTML}{555555}
\definecolor{RuleGray}{HTML}{B8B8B8}
\definecolor{SoftGray}{HTML}{F2F2F0}
\definecolor{PanelGray}{HTML}{FAFAF8}
\definecolor{AccentBlue}{HTML}{1F4E79}
\definecolor{AccentGold}{HTML}{9A6B16}
\definecolor{AccentRed}{HTML}{8F2D2D}
\definecolor{AccentGreen}{HTML}{356B48}

\tikzset{
  paper box/.style={
    draw=Ink,
    rounded corners=2.5pt,
    line width=0.45pt,
    fill=white,
    align=center,
    inner sep=5pt,
    font=\sffamily\small
  },
  paper softbox/.style={
    paper box,
    fill=SoftGray
  },
  paper note/.style={
    draw=RuleGray,
    rounded corners=2pt,
    fill=PanelGray,
    align=left,
    inner sep=5pt,
    font=\sffamily\footnotesize
  },
  paper arrow/.style={
    -{Stealth[length=2.2mm]},
    line width=0.55pt,
    draw=Ink
  },
  paper dashed/.style={
    paper arrow,
    dashed,
    draw=MutedInk
  }
}

\newcommand{\scoreDots}[1]{%
  \ifcase#1\relax
  \or \textbullet\hspace{0.12em}{\color{RuleGray}\textbullet\textbullet}%
  \or \textbullet\textbullet\hspace{0.12em}{\color{RuleGray}\textbullet}%
  \or \textbullet\textbullet\textbullet%
  \fi
}

\newcommand{\figureContextAtGlance}{%
\begin{figure}[H]
\centering
\resizebox{\linewidth}{!}{%
\begin{tikzpicture}[font=\sffamily, node distance=8mm]
  \node[paper softbox, minimum width=0.29\linewidth, minimum height=2.8cm] (objects) {
    \textbf{Knowledge objects}\\[3pt]
    package source\\
    local repo\\
    docs / changelogs\\
    issues / CI logs\\
    memories / skills
  };
  \node[paper softbox, right=14mm of objects, minimum width=0.29\linewidth, minimum height=2.8cm] (routes) {
    \textbf{Routing strategies}\\[3pt]
    vendor\\
    lazy fetch\\
    docs index\\
    semantic search\\
    symbolic search\\
    agentic inspection
  };
  \node[paper softbox, right=14mm of routes, minimum width=0.29\linewidth, minimum height=2.8cm] (constraints) {
    \textbf{Constraints}\\[3pt]
    attention budget\\
    authority\\
    freshness\\
    privacy\\
    latency\\
    trust
  };

  \draw[paper arrow] (objects) -- node[above, font=\sffamily\footnotesize] {select} (routes);
  \draw[paper arrow] (routes) -- node[above, font=\sffamily\footnotesize] {score} (constraints);
  \node[paper box, below=9mm of routes, minimum width=0.78\linewidth] (policy) {
    \textbf{Context router:} choose the smallest sufficient, highest-authority bundle for this task phase.
  };
  \draw[paper arrow] (constraints.south) |- (policy.east);
  \draw[paper arrow] (objects.south) |- (policy.west);
\end{tikzpicture}
}
\caption{Context routing at a glance. The paper's thesis is not that one context source wins; it is that coding agents need an explicit allocation policy over objects, routes, and constraints.}
\label{fig-context-at-glance}
\end{figure}
}

\newcommand{\figureUtilityTradeoff}{%
\begin{figure}[H]
\centering
\begin{tikzpicture}
\begin{axis}[
  width=0.82\linewidth,
  height=5.2cm,
  axis lines=left,
  xlabel={Context supplied},
  ylabel={Expected task utility},
  xmin=0, xmax=10,
  ymin=0, ymax=10,
  xtick=\empty,
  ytick=\empty,
  legend style={draw=none, fill=none, at={(0.5,-0.18)}, anchor=north, legend columns=3, font=\sffamily\footnotesize},
  every axis plot/.append style={line width=1pt},
  label style={font=\sffamily\small},
]
\addplot[AccentBlue, mark=none, smooth] coordinates {(0,1.0) (1,2.4) (2,4.5) (3,6.6) (4,7.7) (5,8.1) (6,8.0) (7,7.3) (8,6.2) (9,4.8) (10,3.0)};
\addlegendentry{net utility}
\addplot[Ink, dashed, mark=none] coordinates {(0,0.6) (1,1.3) (2,2.0) (3,2.7) (4,3.4) (5,4.1) (6,4.9) (7,5.8) (8,6.9) (9,8.0) (10,9.2)};
\addlegendentry{tokens + latency}
\addplot[AccentRed, dotted, mark=none] coordinates {(0,0.3) (1,0.5) (2,0.9) (3,1.3) (4,1.8) (5,2.5) (6,3.4) (7,4.7) (8,6.2) (9,7.9) (10,9.5)};
\addlegendentry{distractor risk}
\node[font=\sffamily\footnotesize, align=center] at (axis cs:4.7,8.8) {smallest\\sufficient\\context};
\draw[-{Stealth[length=2mm]}, AccentBlue] (axis cs:4.8,8.45) -- (axis cs:4.2,7.75);
\end{axis}
\end{tikzpicture}
\caption{The router optimizes expected utility, not raw context size. More context helps until authority and relevance are outweighed by token, latency, staleness, trust, and distractor costs.}
\label{fig-utility-tradeoff}
\end{figure}
}

\newcommand{\figureLongContextDegradation}{%
\begin{figure}[H]
\centering
\begin{tikzpicture}
\begin{groupplot}[
  group style={group size=2 by 1, horizontal sep=1.2cm},
  width=0.46\linewidth,
  height=5.3cm,
  ymin=0,
  ymax=80,
  ymajorgrids=true,
  grid style={draw=RuleGray!45},
  ylabel={Success / score (\%)},
  label style={font=\sffamily\small},
  tick label style={font=\sffamily\footnotesize},
  title style={font=\sffamily\small\bfseries},
  nodes near coords,
  every node near coord/.append style={font=\sffamily\scriptsize, yshift=2pt},
]
\nextgroupplot[
  title={LongSWE-Bench},
  symbolic x coords={32K,256K},
  xtick=data,
]
\addplot[AccentBlue, mark=*, line width=1pt] coordinates {(32K,29) (256K,3)};
\nextgroupplot[
  title={LongCodeQA},
  symbolic x coords={512K,1M},
  xtick=data,
]
\addplot[AccentGold, mark=square*, line width=1pt] coordinates {(512K,70.2) (1M,40.0)};
\end{groupplot}
\end{tikzpicture}
\caption{LongCodeBench examples used in this paper. Claude 3.5 Sonnet drops on LongSWE-Bench from 29\% at 32K to 3\% at 256K; Qwen2.5 drops on LongCodeQA from 70.2\% at 512K to 40.0\% at 1M. The point is calibration, not a universal threshold.}
\label{fig-longcodebench}
\end{figure}
}

\newcommand{\figureRetrievalLandscape}{%
\begin{figure}[H]
\centering
\begin{tikzpicture}[x=1.25cm, y=0.86cm, font=\sffamily\footnotesize]
  \draw[->, line width=0.55pt] (0,0) -- (7.2,0) node[right] {semantic gap handled};
  \draw[->, line width=0.55pt] (0,0) -- (0,5.4) node[above] {determinism / authority};
  \draw[RuleGray] (3.4,0) -- (3.4,5.1);
  \draw[RuleGray] (0,2.55) -- (6.9,2.55);

  \node[paper note, anchor=west] at (0.2,5.0) {exact, local, cheap};
  \node[paper note, anchor=east] at (6.8,5.0) {semantic but still auditable};
  \node[paper note, anchor=west] at (0.2,0.55) {misses intent};
  \node[paper note, anchor=east] at (6.8,0.55) {higher drift / trust cost};

  \node[paper box, minimum width=1.8cm] at (1.2,4.0) {ripgrep\\BM25};
  \node[paper box, minimum width=2.1cm] at (2.7,3.4) {SPLADE\\learned sparse};
  \node[paper box, minimum width=2.0cm] at (4.4,3.1) {hybrid\\RRF};
  \node[paper box, minimum width=1.9cm] at (5.6,2.35) {dense\\embeddings};
  \node[paper box, minimum width=1.9cm] at (5.9,3.9) {ColBERT\\late interaction};
  \node[paper box, minimum width=2.0cm] at (3.4,4.7) {AST / LSP\\SCIP};
  \node[paper box, minimum width=2.0cm] at (6.1,1.3) {agentic\\web/docs};
\end{tikzpicture}
\caption{Semantic and non-semantic search are complements. Exact lexical search is cheap and deterministic; dense and late-interaction search bridge intent gaps; graph and symbol search preserve code authority; hybrid retrieval is often the pragmatic route.}
\label{fig-retrieval-landscape}
\end{figure}
}

\newcommand{\figureFailureModes}{%
\begin{figure}[H]
\centering
\begin{tikzpicture}[font=\sffamily\footnotesize, node distance=3.5mm]
  \foreach \i/\failure/\response in {
    0/Attention dilution/Reduce bundle and place evidence deliberately,
    1/Distractor interference/Prune and prefer version-specific authority,
    2/Representation mismatch/Switch docs to examples or source,
    3/Version skew/Resolve through lockfile provenance,
    4/Retrieval miss/Expand through graph or agentic fallback,
    5/Context staleness/Prefer live local reads or refresh index,
    6/Trust contamination/Quote untrusted context; do not obey it,
    7/Cost blowup/Precompute maps or cap search loops
  }{
    \pgfmathsetmacro{\y}{-\i*0.72}
    \node[paper softbox, minimum width=0.31\linewidth, anchor=west] (f\i) at (0,\y) {\failure};
    \node[paper box, minimum width=0.50\linewidth, anchor=west] (r\i) at (6.0,\y) {\response};
    \draw[paper arrow] (f\i.east) -- (r\i.west);
  }
\end{tikzpicture}
\caption{Wrong context has concrete failure modes. A router is useful only if each failure maps to an intervention: prune, switch representation, refresh, change authority, or stop searching.}
\label{fig-failure-modes}
\end{figure}
}

\newcommand{\heatcell}[1]{%
  \ifcase#1\cellcolor{white}0%
  \or\cellcolor{SoftGray}1%
  \or\cellcolor{RuleGray!45}2%
  \or\cellcolor{Ink!72}\textcolor{white}{3}%
  \fi
}

\newcommand{\figurePhaseHeatmap}{%
\begin{figure}[H]
\centering
\small
\setlength{\tabcolsep}{4pt}
\renewcommand{\arraystretch}{1.12}
\begin{tabular}{lcccccc}
\toprule
\textbf{Phase} & \rotatebox{45}{Rules} & \rotatebox{45}{Repo map} & \rotatebox{45}{Search} & \rotatebox{45}{Files} & \rotatebox{45}{Package src} & \rotatebox{45}{Tests/logs}\\
\midrule
Orient    & \heatcell{3} & \heatcell{3} & \heatcell{1} & \heatcell{1} & \heatcell{0} & \heatcell{0}\\
Locate    & \heatcell{1} & \heatcell{2} & \heatcell{3} & \heatcell{2} & \heatcell{1} & \heatcell{1}\\
Plan      & \heatcell{2} & \heatcell{1} & \heatcell{2} & \heatcell{3} & \heatcell{2} & \heatcell{1}\\
Edit      & \heatcell{1} & \heatcell{0} & \heatcell{1} & \heatcell{3} & \heatcell{2} & \heatcell{1}\\
Verify    & \heatcell{1} & \heatcell{0} & \heatcell{1} & \heatcell{2} & \heatcell{1} & \heatcell{3}\\
Repair    & \heatcell{1} & \heatcell{1} & \heatcell{3} & \heatcell{3} & \heatcell{3} & \heatcell{3}\\
Summarize & \heatcell{2} & \heatcell{2} & \heatcell{0} & \heatcell{2} & \heatcell{0} & \heatcell{1}\\
\bottomrule
\end{tabular}
\caption{Task-phase context heatmap. Darker cells indicate higher expected usefulness. Context routing is an online policy: the right bundle changes as the agent orients, locates, edits, verifies, repairs, and summarizes.}
\label{fig-phase-heatmap}
\end{figure}
}

\newcommand{\figureRouterArchitecture}{%
\begin{figure}[H]
\centering
\resizebox{\linewidth}{!}{%
\begin{tikzpicture}[font=\sffamily\footnotesize, node distance=7mm]
  \node[paper softbox, minimum width=0.86\linewidth] (signals) {
    \textbf{Observable signals:} task type, lockfile, imports, source size, docs quality, model, privacy, failures
  };
  \node[paper box, below=9mm of signals, minimum width=0.30\linewidth, minimum height=1.3cm] (policy) {\textbf{Routing policy}\\utility + constraints};
  \node[paper softbox, left=11mm of policy, minimum width=0.22\linewidth, minimum height=1.1cm] (authority) {authority\\resolver};
  \node[paper softbox, right=11mm of policy, minimum width=0.22\linewidth, minimum height=1.1cm] (trust) {privacy / trust\\gate};
  \node[paper note, below=9mm of policy, minimum width=0.78\linewidth] (routes) {
    routes: vendor source \quad lazy fetch \quad docs index \quad lexical search \quad semantic search \quad graph search \quad skill \quad agentic inspection
  };
  \node[paper box, below left=10mm and 11mm of routes, minimum width=0.30\linewidth] (bundle) {context bundle\\+ tool affordance};
  \node[paper box, below=10mm of routes, minimum width=0.26\linewidth] (lock) {context.lock\\provenance};
  \node[paper box, below right=10mm and 11mm of routes, minimum width=0.30\linewidth] (explain) {explanation\\+ override};

  \draw[paper arrow] (signals) -- (policy);
  \draw[paper arrow] (authority) -- (policy);
  \draw[paper arrow] (trust) -- (policy);
  \draw[paper arrow] (policy) -- (routes);
  \draw[paper arrow] (routes) -- (bundle);
  \draw[paper arrow] (routes) -- (lock);
  \draw[paper arrow] (routes) -- (explain);
\end{tikzpicture}
}
\caption{Context router architecture redrawn as a PDF-native diagram. The router does not merely retrieve content; it resolves authority, enforces trust policy, emits provenance, and explains the route.}
\label{fig-router}
\end{figure}
}

\newcommand{\figureVercelEval}{%
\begin{figure}[H]
\centering
\begin{tikzpicture}
\begin{axis}[
  ybar,
  width=0.78\linewidth,
  height=5.8cm,
  ymin=0,
  ymax=110,
  ylabel={Pass rate (\%)},
  symbolic x coords={Baseline,Skill,Skill+instruction,AGENTS.md},
  xtick=data,
  x tick label style={rotate=18, anchor=east, font=\sffamily\footnotesize},
  ymajorgrids=true,
  grid style={draw=RuleGray!45},
  bar width=18pt,
  nodes near coords,
  every node near coord/.append style={font=\sffamily\scriptsize},
  label style={font=\sffamily\small},
]
\addplot[draw=Ink, fill=SoftGray, postaction={pattern={Lines[angle=45,distance=4pt,line width=0.35pt]}}] coordinates {
  (Baseline,53) (Skill,53) (Skill+instruction,79) (AGENTS.md,100)
};
\end{axis}
\end{tikzpicture}
\caption{Vercel's reported Next.js agent eval results. The lesson for routing is selector and timing: an always-on compact docs index outperformed an on-demand skill in this suite.}
\label{fig-vercel-eval}
\end{figure}
}

\newcommand{\figureEvalLattice}{%
\begin{figure}[H]
\centering
\scriptsize
\setlength{\tabcolsep}{3.2pt}
\renewcommand{\arraystretch}{1.12}
\begin{tabular}{lcccccccc}
\toprule
\textbf{Strategy} & \rotatebox{45}{Unknown lib} & \rotatebox{45}{Framework idiom} & \rotatebox{45}{Migration} & \rotatebox{45}{Stack trace} & \rotatebox{45}{Multi-file} & \rotatebox{45}{Security} & \rotatebox{45}{Perf/debug} & \rotatebox{45}{Test repair}\\
\midrule
No extra context       & -- & -- & -- & -- & o  & o  & -- & o  \\
Docs-only              & +  & o  & +  & -- & -- & +  & -- & o  \\
Always-on docs index   & +  & +  & +  & o  & o  & +  & -- & o  \\
Source vendoring       & o  & +  & o  & +  & o  & o  & +  & +  \\
Lazy source fetch      & o  & o  & o  & +  & o  & o  & +  & +  \\
Repo pack              & o  & o  & -- & o  & +  & -- & o  & o  \\
Semantic index         & +  & o  & o  & o  & +  & o  & o  & o  \\
Structural index       & o  & +  & o  & +  & +  & o  & +  & +  \\
Agentic search         & o  & +  & o  & +  & +  & o  & +  & +  \\
Router                 & +  & +  & +  & +  & +  & +  & +  & +  \\
\bottomrule
\end{tabular}
\caption{Evaluation lattice. Symbols encode expected route relevance: `+' core condition, `o' optional condition, `--' weak or negative-control condition. A router benchmark should vary strategy while holding agent and task fixed.}
\label{fig-eval-lattice}
\end{figure}
}

\newcommand{\figureToolLandscapeCounts}{%
\begin{figure}[H]
\centering
\begin{tikzpicture}
\begin{axis}[
  xbar,
  width=0.84\linewidth,
  height=7.0cm,
  xmin=0,
  xmax=24,
  xlabel={Representative tools counted in snapshot},
  symbolic y coords={Adjacent safety,Multi-repo,Repo packaging,Search/indexing,Embeddings/RAG,MCP servers,IDE context engines},
  ytick=data,
  y dir=reverse,
  nodes near coords,
  every node near coord/.append style={font=\sffamily\scriptsize},
  tick label style={font=\sffamily\footnotesize},
  label style={font=\sffamily\small},
  xmajorgrids=true,
  grid style={draw=RuleGray!45},
  bar width=9pt,
]
\addplot[draw=Ink, fill=SoftGray, postaction={pattern={Lines[angle=0,distance=4pt,line width=0.35pt]}}] coordinates {
  (7,Adjacent safety)
  (7,Multi-repo)
  (7,Repo packaging)
  (14,Search/indexing)
  (16,Embeddings/RAG)
  (20,MCP servers)
  (21,IDE context engines)
};
\end{axis}
\end{tikzpicture}
\caption{Tool-landscape fragmentation in the May 2026 snapshot. The distribution itself motivates routing: tools cluster around different context surfaces rather than one universal context strategy.}
\label{fig-tool-counts}
\end{figure}
}



\title{Context Routing for AI Coding Agents}
\usepackage{etoolbox}
\makeatletter
\providecommand{\subtitle}[1]{% add subtitle to \maketitle
  \apptocmd{\@title}{\par {\large #1 \par}}{}{}
}
\makeatother
\subtitle{A Research Agenda for Context Allocation Under Attention,
Authority, Freshness, Privacy, and Cost Constraints}
\author{Gunther Brunner}
\date{2026-05-13}
\begin{document}
\maketitle
\begin{abstract}
AI coding agents increasingly fail not because relevant knowledge is
unavailable, but because it is poorly allocated: too much, too stale,
retrieved from the wrong authority, positioned where the model attends
weakly, or exposed through the wrong trust boundary. Context windows are
storage budgets, not attention guarantees. A million-token window lets
an agent carry more material; it does not imply that the model will
retrieve, prioritize, and apply the right evidence under task pressure.

This paper surveys academic evidence on long-context degradation,
retrieval-augmented code generation, repository-level retrieval,
structural code search, and distractor sensitivity, alongside
practitioner designs from source vendoring, agentic search, code
indexing, repo maps, managed context engines, and progressive-disclosure
skills. The synthesis is simple: no single context strategy dominates
across model, task, repository, package, privacy, freshness, and latency
conditions. The correct unit of design is a context-routing policy.

We therefore propose the context router: an explicit policy layer that
selects, per task and knowledge object, among vendored source, lazy
source fetch, documentation retrieval, semantic or structural search,
packaged snapshots, skills, and agent-mediated inspection. We formalize
the routing problem, separate the design space into orthogonal axes,
identify failure modes and policy signals, define source authority and
provenance requirements, outline a router-level evaluation agenda, and
map current standards onto the abstraction. The paper is motivated by
work on ingraft, but its primary claim is broader: context routing
should become an inspectable, overrideable, standards-based layer rather
than an opaque feature hidden inside individual coding agents.
\end{abstract}

\renewcommand*\contentsname{Table of contents}
{
\hypersetup{linkcolor=black}
\setcounter{tocdepth}{3}
\tableofcontents
}

\section{Introduction}\label{sec-intro}

AI coding agents are now useful enough that their failures matter.
Developers ask them to patch production repositories, migrate
frameworks, inspect failing tests, answer questions about large
codebases, and integrate unfamiliar libraries. In these settings, the
bottleneck is rarely ``no relevant information exists.'' The bottleneck
is context allocation: which information should enter the agent's
working set, in which representation, through which trust boundary, at
what point in the loop, and under what budget.

The central claim of this paper is:

\begin{quote}
No single context strategy dominates. The correct unit of design is a
context-routing policy.
\end{quote}

\figureContextAtGlance

That claim is deliberately broader than source vendoring or any single
product architecture. A coding agent needs many kinds of knowledge:
package source, repository source, local conventions, issue threads,
stack traces, failing test logs, API docs, examples, migration guides,
previous decisions, generated summaries, dependency lockfiles, CI logs,
design docs, architecture diagrams, code-search indexes, human
instructions, and memory across sessions. A package is only one context
object.

The more general abstraction is:

\begin{quote}
A context router selects a representation, retrieval mechanism, trust
boundary, and loading schedule for each knowledge object relevant to a
coding task.
\end{quote}

In that framing, ingraft is not the proof of the thesis. It is a case
study: ingraft implements the package-knowledge slice of a larger
context-allocation problem. For third-party package knowledge, the
router may choose vendored source, lazy source fetch, docs retrieval,
semantic search, packaged snapshots, skills, or agentic inspection. For
other knowledge objects, the same routing logic applies to issue
trackers, project docs, CI logs, API specs, or local architecture notes.

The immediate product trigger was Maxwell Brown's Effect post on using
git subtrees for coding-agent context (Brown 2026). Brown's short
version was ``Coding agents rely on patterns and examples, not
descriptions'' and ``give your agent real library code to learn from.''
That claim is intentionally narrower than this paper's thesis. It does
not settle the general problem; it exposes it: if subtree is the right
route for Effect, what policy decides when subtree is right, when docs
are better, and when neither should be loaded?

This paper contributes five things:

\begin{enumerate}
\def\labelenumi{\arabic{enumi}.}
\tightlist
\item
  A problem formulation for coding-agent context allocation.
\item
  An evidence synthesis that distinguishes peer-reviewed research,
  preprints, vendor engineering posts, practitioner experience,
  standards, and the author's interpretation.
\item
  A multidimensional design space for context strategies: placement,
  representation, selector, timing, and constraints.
\item
  A context-router design model covering policy signals, source
  authority, provenance, trust classes, task-phase routing, and context
  lockfiles.
\item
  A router-level evaluation agenda with ablations and falsification
  tests.
\end{enumerate}

The intended result is not a larger landscape memo. It is a sharper
argument: context engineering for coding agents is routing under
attention, authority, freshness, privacy, and cost constraints. Existing
tools each solve one region of that design space. A router is the
abstraction that makes those partial solutions composable.

\section{Method and Scope}\label{sec-method}

This paper is a dated synthesis, not an exhaustive census. The review
covers materials visible through May 13, 2026, with most tool and
standards claims checked against public documentation, repositories,
specifications, or vendor posts available during March-May 2026.

\subsection{Search and inclusion criteria}\label{sec-method-search}

Sources searched included arXiv, ACL Anthology, ACM/IEEE proceedings
where available, OpenReview, GitHub, vendor documentation, engineering
blogs, standards websites, tool registries, Hacker News, practitioner
podcasts, and industry surveys.

Papers were included when they addressed at least one of these
questions:

\begin{itemize}
\tightlist
\item
  long-context evaluation for language models or coding models;
\item
  repository-level code generation, completion, repair, or localization;
\item
  retrieval-augmented code generation;
\item
  semantic, symbolic, graph, or AST-based code retrieval;
\item
  agent scaffolding, context management, compression, or distractor
  behavior.
\end{itemize}

Tools were included when they affected how coding agents obtain
repository, library, documentation, organizational, or procedural
context. The appendix is a representative May 2026 snapshot, not a
complete inventory. Tool counts and product features are volatile.

Excluded materials include pure model benchmarks with no context
intervention, general-purpose RAG systems with no code or agent
relevance, and abandoned tools unless they were historically
influential.

\subsection{Evidence hierarchy}\label{sec-evidence-hierarchy}

The paper uses heterogeneous evidence. To keep claims calibrated,
Table~\ref{tbl-evidence} separates the role of each evidence class.

\begin{longtable}[]{@{}
  >{\raggedright\arraybackslash}p{(\linewidth - 4\tabcolsep) * \real{0.3333}}
  >{\raggedright\arraybackslash}p{(\linewidth - 4\tabcolsep) * \real{0.3333}}
  >{\raggedright\arraybackslash}p{(\linewidth - 4\tabcolsep) * \real{0.3333}}@{}}
\caption{Evidence classes and their role in the
argument.}\label{tbl-evidence}\tabularnewline
\toprule\noalign{}
\begin{minipage}[b]{\linewidth}\raggedright
Evidence class
\end{minipage} & \begin{minipage}[b]{\linewidth}\raggedright
Examples
\end{minipage} & \begin{minipage}[b]{\linewidth}\raggedright
How it is used
\end{minipage} \\
\midrule\noalign{}
\endfirsthead
\toprule\noalign{}
\begin{minipage}[b]{\linewidth}\raggedright
Evidence class
\end{minipage} & \begin{minipage}[b]{\linewidth}\raggedright
Examples
\end{minipage} & \begin{minipage}[b]{\linewidth}\raggedright
How it is used
\end{minipage} \\
\midrule\noalign{}
\endhead
\bottomrule\noalign{}
\endlastfoot
Peer-reviewed empirical papers & Lost in the Middle, RULER, RepoBench,
SWE-bench, GraphCoder & Mechanisms and benchmark results with
comparatively strong evidential weight \\
Preprints and technical reports & LongCodeBench, cAST, Context Rot,
newer repo-retrieval papers & Emerging findings, cited with venue
caveats \\
Vendor engineering posts & Anthropic, Cursor, Cognition, Augment,
Sourcegraph & Evidence of deployed design patterns and product
constraints, not independent proof \\
Practitioner essays and podcasts & Effect, Cline, aider, HumanLayer,
Ptacek, Ronacher & Operating experience and design disagreement \\
Tool documentation and standards & MCP, Agent Skills, AGENTS.md,
llms.txt, SCIP & Implementation facts and interoperability surfaces \\
This paper's synthesis & router thesis, policy model, evaluation design
& Interpretation that should be tested, not treated as direct empirical
result \\
\end{longtable}

Two terms are used carefully:

\begin{itemize}
\tightlist
\item
  \textbf{Academic and technical works} means peer-reviewed papers,
  workshop papers, arXiv preprints, benchmark reports, and technical
  reports. It does not imply every cited work is peer-reviewed.
\item
  \textbf{Representative tool snapshot} means a dated inventory of
  visible major tools. It is not complete and will age.
\end{itemize}

Practitioner evidence shows convergence around a shared problem, not
consensus on the solution. The field agrees that context selection
matters; it disagrees on where selection should happen.

\section{Definitions and Problem Formulation}\label{sec-formal}

\subsection{Core definitions}\label{sec-definitions}

\textbf{Context} is any information made available to an agent during a
task. It includes prompt tokens, files the agent can read, tool results,
memory, indexes, skills, instructions, logs, and external service
responses.

\textbf{Context object} is a knowledge unit that might be relevant: a
package, file, repository, issue, stack trace, lockfile entry, API
document, design note, generated summary, or skill.

\textbf{Context bundle} is the concrete material exposed to the agent
for a task: selected files, docs, snippets, index hits, rules,
provenance notes, and tool affordances.

\textbf{Routing strategy} is a way of exposing a context object: vendor
it into the workspace, fetch it lazily, retrieve docs, query a graph
index, pack a repo, activate a skill, or let an agent inspect through
shell tools.

\textbf{Context router} is a policy
\texttt{\textbackslash{}pi(t,\ m,\ r,\ k)\ -\textgreater{}\ s}: given a
task \texttt{t}, model or agent capability \texttt{m}, repository state
\texttt{r}, and knowledge object \texttt{k}, select a routing strategy
\texttt{s}.

\textbf{Source authority} describes which source should win for a
specific question: the lockfile for installed version, official docs for
public API intent, source and tests for behavior, local code for project
convention, and human instructions for constraints.

\textbf{Context provenance} records where a bundle came from, which
version it represents, when it was fetched, how trustworthy it is, and
what it is intended to be used for.

\subsection{Utility model}\label{sec-utility}

A lightweight model is enough to make the tradeoff explicit. For a
coding task \texttt{t}, model or agent \texttt{m}, repository state
\texttt{r}, knowledge object \texttt{k}, and candidate strategy
\texttt{s}, the router wants high expected success and low total cost:

\begin{align*}
U(s \mid t,m,r,k) ={}&
P(\mathrm{success} \mid t,m,r,C_s) \\
&- \lambda_1\mathrm{tokens}(C_s)
- \lambda_2\mathrm{latency}(s)
- \lambda_3\mathrm{privacyRisk}(s) \\
&- \lambda_4\mathrm{staleness}(s)
- \lambda_5\mathrm{distractorRisk}(C_s)
- \lambda_6\mathrm{trustRisk}(s)
\end{align*}

The context router chooses the strategy with the highest expected
utility under observable signals. The exact weights are product and
organization specific, but the terms make clear why ``more context'' is
not the objective.

\figureUtilityTradeoff

The operational principle is \textbf{smallest sufficient context}:

\begin{quote}
Context is sufficient when it lets the agent identify the relevant
authority, produce a correct change, and verify it without introducing
materially misleading distractors.
\end{quote}

Smallest does not mean fewest tokens. It means lowest total context cost
across tokens, latency, privacy, staleness, trust, and distractor risk.
Sometimes the lowest-cost context is a small doc page. Sometimes it is a
large but version-matched source tree.

\subsection{Policy signals}\label{sec-policy-signals}

The router cannot observe expected success directly, so it uses signals.
The signals in Table~\ref{tbl-signals} are enough to turn ``router''
from metaphor into a policy surface.

\begin{longtable}[]{@{}
  >{\raggedright\arraybackslash}p{(\linewidth - 2\tabcolsep) * \real{0.5000}}
  >{\raggedright\arraybackslash}p{(\linewidth - 2\tabcolsep) * \real{0.5000}}@{}}
\caption{Observable signals for a context-routing
policy.}\label{tbl-signals}\tabularnewline
\toprule\noalign{}
\begin{minipage}[b]{\linewidth}\raggedright
Signal
\end{minipage} & \begin{minipage}[b]{\linewidth}\raggedright
Why it matters
\end{minipage} \\
\midrule\noalign{}
\endfirsthead
\toprule\noalign{}
\begin{minipage}[b]{\linewidth}\raggedright
Signal
\end{minipage} & \begin{minipage}[b]{\linewidth}\raggedright
Why it matters
\end{minipage} \\
\midrule\noalign{}
\endhead
\bottomrule\noalign{}
\endlastfoot
Package appears in lockfile & Fixes version-specific context and avoids
latest-doc drift \\
Import frequency & Estimates package centrality and likely value of
deeper source \\
Package source size & Affects vendoring, packing, indexing, and
distractor cost \\
Docs availability and quality & Makes docs-first viable for public API
tasks \\
API volatility & Increases value of fresh docs, changelogs, or source \\
Local modifications & Penalizes stale indexes and remote snapshots \\
Model context window and measured effective length & Bounds what can be
reliably carried \\
Agent tool capability & Determines whether agentic search is viable \\
Privacy policy & Rules out cloud retrieval or remote indexing \\
Task type & Bugfix, migration, greenfield integration, test repair,
security review, or refactor \\
Failure history & Repeated failures justify deeper context or a strategy
switch \\
Repository size and churn & Determines index/search break-even \\
Language ecosystem support & Determines whether AST, LSP, or SCIP
retrieval is available \\
\end{longtable}

\section{Why Context Allocation Is Hard}\label{sec-hard}

\subsection{Context windows are not attention
guarantees}\label{sec-attention}

Long windows increase storage, not reliable computation over every
token. N. F. Liu et al. (2024) established the now-canonical positional
failure: models perform best when relevant information is near the
beginning or end of the context and degrade when evidence is buried in
the middle. {Hsieh et al.} (2024) showed that advertised context length
can overstate effective context length, especially on multi-hop
synthetic tasks.

LongCodeBench extends the problem to realistic code comprehension and
repair (Rando et al. 2025). It evaluates LongCodeQA and LongSWE-Bench
across 32K, 64K, 128K, 256K, 512K, and 1M context brackets where
supported. The precise headline is:

\begin{quote}
On LongSWE-Bench, Claude 3.5 Sonnet solved 29\% of issues at 32K context
and 3\% at 256K; it was not a 1M-token Claude result. On LongCodeQA,
Qwen2.5 peaked at 70.2\% at 512K and fell to 40.0\% at 1M.
\end{quote}

\figureLongContextDegradation

The newer context-rot literature sharpens the same point. Chroma's
technical report evaluates 18 models and isolates cases where increasing
input length alone makes outputs less reliable, including semantic
needle variants, distractors, LongMemEval comparisons, and repeated-word
copying tasks (Hong, Troynikov, and Huber 2025). NoLiMa extends
needle-in-a-haystack beyond literal matching and reports severe drops
when questions and needles have little lexical overlap, even when short
contexts are nearly solved (Modarressi et al. 2025). AbsenceBench adds a
different failure class: models that can retrieve present evidence may
still fail to identify missing information, including omitted lines in
GitHub pull-request diffs (Fu et al. 2025).

The defensible conclusion is not ``1M always fails.'' The conclusion is
that advertised context length is an upper bound, not an operating
target. Routing budgets should be calibrated per model, task class, and
representation, ideally by local evals or conservative defaults.

\subsection{Retrieval can help or hurt}\label{sec-retrieval-hard}

{Z. Z. Wang et al.} (2025) asks directly whether retrieval can augment
code generation. Its answer is conditional: high-quality retrieved
contexts improve some settings, especially weaker models and basic
programming tasks, while retrievers still struggle when lexical overlap
is limited and generators do not always integrate extra context
effectively. Asai et al. (2024) reaches a similar architectural lesson
from a broader RAG setting: indiscriminate fixed retrieval can reduce
versatility and produce unhelpful generations; retrieval should be
adaptive.

The router should not ask ``RAG or not?'' It should ask which retrieval
substrate has the lowest expected miss, distractor, staleness, and
latency cost for this task.

That means semantic and non-semantic search should be treated as
different instruments, not rival slogans. BM25 remains a strong lexical
baseline because exact terms, identifiers, and error strings often are
the task (Robertson and Zaragoza 2009). Learned sparse retrieval such as
SPLADE keeps lexical auditability while adding expansion terms (Formal,
Piwowarski, and Clinchant 2021). Dense retrieval and text embeddings are
useful when intent does not share surface form with the target, but
benchmarks such as BEIR and MTEB show that retrieval quality is highly
task- and distribution-dependent (Thakur et al. 2021; Muennighoff et al.
2023). CodeSearchNet made the natural-language-to-code gap explicit for
semantic code search (Husain et al. 2019), while GraphCodeBERT shows why
code structure and data flow matter beyond flat token embeddings (Guo et
al. 2021). Late-interaction retrievers such as ColBERT occupy a useful
middle ground: they retain token-level matching signals while allowing
document representations to be precomputed (Khattab and Zaharia 2020).
ColBERTv2 and PLAID make the deployment point stronger by reducing
late-interaction storage and serving cost without giving up the core
matching advantage (Santhanam et al. 2022a; Santhanam et al. 2022b).
Mixedbread's mxbai work is relevant as practitioner evidence for both
dense embeddings and ColBERT-style reranking, and mgrep is a concrete
coding-agent product claim that semantic search can reduce token usage
relative to repeated grep loops while still preserving grep as the exact
match fallback (Lee et al. 2024a; Lee et al. 2024b; Mixedbread 2026).
Multi-function embedding systems such as BGE-M3 push in the same
direction by supporting dense, sparse, and multi-vector retrieval modes
inside one model family (Chen et al. 2024).

Recent code and technical-document retrieval benchmarks make the same
lesson more concrete for software agents. FreshStack argues for
benchmarks built from recent technical documentation, where off-the-shelf
retrievers can lag far behind oracle context on niche stacks (Thakur et
al. 2025). COIR separates code retrieval into multiple task families
rather than treating ``code search'' as one problem (Li et al. 2024).
CoSIL frames issue localization as graph-guided repository search rather
than flat chunk retrieval (Jiang et al. 2025). For routing, these papers
matter because they move the selector question from ``semantic search or
grep?'' to ``which retrieval substrate matches this task, authority, and
failure mode?''

\figureRetrievalLandscape

\begin{longtable}[]{@{}
  >{\raggedright\arraybackslash}p{(\linewidth - 2\tabcolsep) * \real{0.5000}}
  >{\raggedright\arraybackslash}p{(\linewidth - 2\tabcolsep) * \real{0.5000}}@{}}
\caption{Retrieval substrates and their common failure
modes.}\label{tbl-retrieval-types}\tabularnewline
\toprule\noalign{}
\begin{minipage}[b]{\linewidth}\raggedright
Retrieval type
\end{minipage} & \begin{minipage}[b]{\linewidth}\raggedright
Typical failure mode
\end{minipage} \\
\midrule\noalign{}
\endfirsthead
\toprule\noalign{}
\begin{minipage}[b]{\linewidth}\raggedright
Retrieval type
\end{minipage} & \begin{minipage}[b]{\linewidth}\raggedright
Typical failure mode
\end{minipage} \\
\midrule\noalign{}
\endhead
\bottomrule\noalign{}
\endlastfoot
Lexical grep & Misses semantic matches and renamed concepts \\
Semantic vector search & Retrieves plausible but wrong chunks \\
Graph retrieval & Depends on graph quality and language support \\
AST chunking & Can miss runtime behavior and dynamic dispatch \\
Docs retrieval & Docs may be stale, incomplete, or for the wrong
version \\
Agentic search & Slow, non-deterministic, and tool-call heavy \\
Hybrid search & More accurate but harder to debug and reproduce \\
\end{longtable}

\subsection{Structure matters}\label{sec-structure}

Repository-level work consistently shows that flat text is a weak
abstraction for code. F. Zhang et al. (2023), Shrivastava et al. (2023),
Ding et al. (2024), {Ding et al.} (2023), {Liang et al.} (2024), {Phan
et al.} (2025), and {W. Liu et al.} (2024) all point toward the same
lesson: imports, definitions, call relationships, type information,
symbols, and file structure provide routing signals that raw chunks do
not. {Y. (Sherry). Zhang et al.} (2025) adds a newer preprint result for
AST-aware chunking.

The careful implication is not ``AST plus graph must always be the
default.'' It is: for repository-scale retrieval, the prior should favor
structural retrieval whenever the language ecosystem and latency budget
permit it.

\subsection{Wrong context has concrete harms}\label{sec-wrong-context}

Context quality fails through recognizable mechanisms.
Table~\ref{tbl-failure-modes} maps each failure mode to a router
response.

\begin{longtable}[]{@{}
  >{\raggedright\arraybackslash}p{(\linewidth - 4\tabcolsep) * \real{0.3333}}
  >{\raggedright\arraybackslash}p{(\linewidth - 4\tabcolsep) * \real{0.3333}}
  >{\raggedright\arraybackslash}p{(\linewidth - 4\tabcolsep) * \real{0.3333}}@{}}
\caption{Failure modes that motivate
routing.}\label{tbl-failure-modes}\tabularnewline
\toprule\noalign{}
\begin{minipage}[b]{\linewidth}\raggedright
Failure mode
\end{minipage} & \begin{minipage}[b]{\linewidth}\raggedright
Description
\end{minipage} & \begin{minipage}[b]{\linewidth}\raggedright
Router response
\end{minipage} \\
\midrule\noalign{}
\endfirsthead
\toprule\noalign{}
\begin{minipage}[b]{\linewidth}\raggedright
Failure mode
\end{minipage} & \begin{minipage}[b]{\linewidth}\raggedright
Description
\end{minipage} & \begin{minipage}[b]{\linewidth}\raggedright
Router response
\end{minipage} \\
\midrule\noalign{}
\endhead
\bottomrule\noalign{}
\endlastfoot
Attention dilution & Relevant tokens are present but weakly attended &
Reduce bundle size; place high-value context early or late \\
Distractor interference & Plausible irrelevant context changes the
reasoning path & Prune aggressively; prefer version-specific
authority \\
Representation mismatch & Docs explain concepts but not usage patterns &
Fetch examples, source, or tests \\
Version skew & Retrieved docs/source do not match installed package &
Resolve through lockfile and provenance \\
Retrieval miss & Search fails to fetch the necessary file or symbol &
Expand via graph, AST, LSP, or agentic search \\
Context staleness & Index reflects old repo state & Prefer live grep or
incremental index \\
Trust contamination & External docs contain malicious or irrelevant
instructions & Quote, sanitize, label provenance, sandbox tools \\
Cost blowup & Agent spends many turns searching & Precompute repo maps
or local indexes \\
\end{longtable}

\figureFailureModes

Examples are easy to find in everyday agent work: the agent uses latest
docs for an older installed package, imports a private symbol after
reading source, acts on a StackOverflow answer for the wrong major
version, searches a stale index while local uncommitted changes matter,
follows instructions embedded in third-party docs, or receives a whole
repo pack and overlooks the relevant file in the middle. These are not
model failures alone. They are allocation failures.

\subsection{Context changes during the agent loop}\label{sec-phases}

Context routing is online, not a one-time preprocessing step. Each phase
of an agent loop wants different material.

\begin{longtable}[]{@{}
  >{\raggedright\arraybackslash}p{(\linewidth - 2\tabcolsep) * \real{0.5000}}
  >{\raggedright\arraybackslash}p{(\linewidth - 2\tabcolsep) * \real{0.5000}}@{}}
\caption{Task-phase routing.}\label{tbl-phases}\tabularnewline
\toprule\noalign{}
\begin{minipage}[b]{\linewidth}\raggedright
Phase
\end{minipage} & \begin{minipage}[b]{\linewidth}\raggedright
Useful context
\end{minipage} \\
\midrule\noalign{}
\endfirsthead
\toprule\noalign{}
\begin{minipage}[b]{\linewidth}\raggedright
Phase
\end{minipage} & \begin{minipage}[b]{\linewidth}\raggedright
Useful context
\end{minipage} \\
\midrule\noalign{}
\endhead
\bottomrule\noalign{}
\endlastfoot
Orient & AGENTS.md, repo map, architecture summary, task brief \\
Locate & Code search, symbols, imports, stack traces, issue
references \\
Plan & Relevant local files, docs, examples, package conventions \\
Edit & Exact files, types, tests, local wrappers, public API
contracts \\
Verify & Test logs, CI output, type errors, runtime traces \\
Repair & Failing trace, package source, issue history, migration docs \\
Summarize & Decision log, changed files, unresolved risks, provenance \\
\end{longtable}

\figurePhaseHeatmap

This aligns with Anthropic's context-engineering framing: curation and
maintenance of useful tokens happens throughout inference, not only at
prompt construction time (Anthropic 2025a). Emerging OSS-focused context
engineering work makes the same point at repository scale: agent
configuration files, guidance documents, and tool manifests vary widely,
which means context policy itself becomes part of the software
architecture rather than an afterthought (Mohsenimofidi et al. 2025).

\section{Design Space}\label{sec-design-space}

The earlier ``seven primitives'' vocabulary is useful as shorthand, but
it mixes dimensions. Vendoring, indexing, retrieving, packing, skills,
annotations, and agentic reads are common bundles across several
independent axes. A more useful taxonomy separates the axes first.

\subsection{Axis 1: placement}\label{sec-axis-placement}

\begin{longtable}[]{@{}
  >{\raggedright\arraybackslash}p{(\linewidth - 2\tabcolsep) * \real{0.5000}}
  >{\raggedright\arraybackslash}p{(\linewidth - 2\tabcolsep) * \real{0.5000}}@{}}
\caption{Where context lives.}\label{tbl-placement}\tabularnewline
\toprule\noalign{}
\begin{minipage}[b]{\linewidth}\raggedright
Placement
\end{minipage} & \begin{minipage}[b]{\linewidth}\raggedright
Examples
\end{minipage} \\
\midrule\noalign{}
\endfirsthead
\toprule\noalign{}
\begin{minipage}[b]{\linewidth}\raggedright
Placement
\end{minipage} & \begin{minipage}[b]{\linewidth}\raggedright
Examples
\end{minipage} \\
\midrule\noalign{}
\endhead
\bottomrule\noalign{}
\endlastfoot
In prompt & Pasted file, repo pack, selected snippets \\
On disk in workspace & Vendored source, cloned dependency, exported
docs \\
Local cache & OpenSrc-style package cache, tool-managed snapshots \\
Local index & Sourcebot, Serena, local vector DB, SCIP/LSP graph \\
Remote service & Context7, Augment, Glean, hosted semantic search \\
Model or tool memory & Skills, summaries, persistent notes, agent
memory \\
\end{longtable}

\subsection{Axis 2: representation}\label{sec-axis-representation}

\begin{longtable}[]{@{}ll@{}}
\caption{How context is
represented.}\label{tbl-representation}\tabularnewline
\toprule\noalign{}
Representation & Examples \\
\midrule\noalign{}
\endfirsthead
\toprule\noalign{}
Representation & Examples \\
\midrule\noalign{}
\endhead
\bottomrule\noalign{}
\endlastfoot
Raw source & Vendored repository, fetched package source \\
Human docs & API docs, README, migration guides, changelogs \\
Symbol graph & SCIP, LSP, AST, call/reference graph \\
Embedding chunks & Vector search over files, docs, or issues \\
Summaries & Repo maps, architecture notes, skills, memory \\
Executable tools & MCP tools, scripts, local CLI commands \\
Instructions & AGENTS.md, CLAUDE.md, project rules \\
\end{longtable}

\subsection{Axis 3: selector}\label{sec-axis-selector}

\begin{longtable}[]{@{}ll@{}}
\caption{Who or what selects
context.}\label{tbl-selector}\tabularnewline
\toprule\noalign{}
Selector & Examples \\
\midrule\noalign{}
\endfirsthead
\toprule\noalign{}
Selector & Examples \\
\midrule\noalign{}
\endhead
\bottomrule\noalign{}
\endlastfoot
Human & \texttt{@file}, explicit package, manual doc link \\
Static policy & allowlist routes \texttt{effect} to vendor \\
Heuristic router & import count, package size, privacy class \\
Learned retriever & embedding search, reranker \\
Graph retriever & symbol, dependency, call/reference expansion \\
Agentic search & \texttt{grep}/read/bash loop \\
Oracle/evaluator & benchmark-only upper bound \\
\end{longtable}

\subsection{Axis 4: timing}\label{sec-axis-timing}

\begin{longtable}[]{@{}ll@{}}
\caption{When context is loaded.}\label{tbl-timing}\tabularnewline
\toprule\noalign{}
Timing & Examples \\
\midrule\noalign{}
\endfirsthead
\toprule\noalign{}
Timing & Examples \\
\midrule\noalign{}
\endhead
\bottomrule\noalign{}
\endlastfoot
Always loaded & AGENTS.md, compact repo guidance \\
Session start & Repo map, recent decision log \\
Task start & Package docs, relevant tests, task issue \\
Just in time & Agent calls search, router fetches source \\
After failure & Fetch package source after stack trace points there \\
Compacted later & Summarize stale context for continuation \\
\end{longtable}

\subsection{Axis 5: dominant constraints}\label{sec-axis-constraints}

\begin{longtable}[]{@{}ll@{}}
\caption{Constraints that dominate route
choice.}\label{tbl-constraints}\tabularnewline
\toprule\noalign{}
Constraint & Examples \\
\midrule\noalign{}
\endfirsthead
\toprule\noalign{}
Constraint & Examples \\
\midrule\noalign{}
\endhead
\bottomrule\noalign{}
\endlastfoot
Token budget & Long-context degradation, prompt cost \\
Latency & Agentic search and remote retrieval delays \\
Privacy & Local-only code, no cloud index, no remote docs \\
Freshness & Local changes, docs drift, moving issues \\
Trust & MCP/tool/skill supply chain, prompt injection \\
Determinism & Reproducible context versus live retrieval \\
Version specificity & Lockfile-resolved package versions \\
\end{longtable}

The old primitives can now be expressed more precisely:

\begin{itemize}
\tightlist
\item
  \textbf{Vendoring} = workspace placement + raw source representation +
  static or heuristic selection + task-time loading + strong
  reproducibility + local privacy.
\item
  \textbf{Indexing} = local or remote placement + symbol/embedding
  representation

  \begin{itemize}
  \tightlist
  \item
    learned or graph selector + low-latency query + staleness
    management.
  \end{itemize}
\item
  \textbf{Skills} = memory placement + instruction/tool representation +
  metadata-based selector + progressive disclosure + high token
  efficiency.
\item
  \textbf{Repo pack} = prompt placement + raw or summarized
  representation + human or static selector + one-shot timing + weak
  freshness after creation.
\end{itemize}

For implementation, these axes become a route descriptor rather than a
prose taxonomy.

\begin{longtable}[]{@{}
  >{\raggedright\arraybackslash}p{(\linewidth - 4\tabcolsep) * \real{0.3333}}
  >{\raggedright\arraybackslash}p{(\linewidth - 4\tabcolsep) * \real{0.3333}}
  >{\raggedright\arraybackslash}p{(\linewidth - 4\tabcolsep) * \real{0.3333}}@{}}
\caption{Route-descriptor axes for a context
object.}\label{tbl-route-descriptor}\tabularnewline
\toprule\noalign{}
\begin{minipage}[b]{\linewidth}\raggedright
Axis
\end{minipage} & \begin{minipage}[b]{\linewidth}\raggedright
Question
\end{minipage} & \begin{minipage}[b]{\linewidth}\raggedright
Example values
\end{minipage} \\
\midrule\noalign{}
\endfirsthead
\toprule\noalign{}
\begin{minipage}[b]{\linewidth}\raggedright
Axis
\end{minipage} & \begin{minipage}[b]{\linewidth}\raggedright
Question
\end{minipage} & \begin{minipage}[b]{\linewidth}\raggedright
Example values
\end{minipage} \\
\midrule\noalign{}
\endhead
\bottomrule\noalign{}
\endlastfoot
Object kind & What is being routed? & Package source, repo subtree, doc
page, API spec, ticket, skill, live endpoint \\
Materialization & Where does it live for the task? & Committed, cached,
packed, mounted, remote-only \\
Representation & In what form is it exposed? & Raw source, structural
summary, prose doc, metadata envelope, search hit \\
Access path & How does the agent reach it? & Direct read, local grep,
local index, remote retrieval API, MCP tool \\
Freshness mode & How stable is it? & Pinned snapshot, refreshable cache,
live mutable \\
Authority/trust & How much should the agent rely on it? & Canonical,
upstream authoritative, derived, advisory, unverified \\
Privacy scope & What boundary constrains it? & Repo-safe, local-only,
restricted, remote-call-required \\
\end{longtable}

\section{Practitioner Strategies as Operating
Zones}\label{sec-practitioner}

The practitioner landscape is not a consensus. It is structured
disagreement. Each camp has a correct operating zone and a failure mode.

\begin{longtable}[]{@{}
  >{\raggedright\arraybackslash}p{(\linewidth - 10\tabcolsep) * \real{0.1667}}
  >{\raggedright\arraybackslash}p{(\linewidth - 10\tabcolsep) * \real{0.1667}}
  >{\raggedright\arraybackslash}p{(\linewidth - 10\tabcolsep) * \real{0.1667}}
  >{\raggedright\arraybackslash}p{(\linewidth - 10\tabcolsep) * \real{0.1667}}
  >{\raggedright\arraybackslash}p{(\linewidth - 10\tabcolsep) * \real{0.1667}}
  >{\raggedright\arraybackslash}p{(\linewidth - 10\tabcolsep) * \real{0.1667}}@{}}
\caption{Practitioner camps, assumptions, failure modes, and routing
implications.}\label{tbl-camps}\tabularnewline
\toprule\noalign{}
\begin{minipage}[b]{\linewidth}\raggedright
Camp
\end{minipage} & \begin{minipage}[b]{\linewidth}\raggedright
Claim
\end{minipage} & \begin{minipage}[b]{\linewidth}\raggedright
Evidence
\end{minipage} & \begin{minipage}[b]{\linewidth}\raggedright
Hidden assumption
\end{minipage} & \begin{minipage}[b]{\linewidth}\raggedright
Failure mode
\end{minipage} & \begin{minipage}[b]{\linewidth}\raggedright
Router implication
\end{minipage} \\
\midrule\noalign{}
\endfirsthead
\toprule\noalign{}
\begin{minipage}[b]{\linewidth}\raggedright
Camp
\end{minipage} & \begin{minipage}[b]{\linewidth}\raggedright
Claim
\end{minipage} & \begin{minipage}[b]{\linewidth}\raggedright
Evidence
\end{minipage} & \begin{minipage}[b]{\linewidth}\raggedright
Hidden assumption
\end{minipage} & \begin{minipage}[b]{\linewidth}\raggedright
Failure mode
\end{minipage} & \begin{minipage}[b]{\linewidth}\raggedright
Router implication
\end{minipage} \\
\midrule\noalign{}
\endhead
\bottomrule\noalign{}
\endlastfoot
Source vendoring & Agents learn patterns better from source than prose &
Effect subtree argument (Brown 2026) & Source is available, navigable,
version-matched, and not huge & Internal APIs, large repos, poor
comments, private access & Vendor high-centrality packages with
manageable source and high idiom value \\
Agentic search & Let the model inspect live files instead of maintaining
an index & Claude Code and Cline practice (S. (swyx). Wang and Trotta
2025; Baumann 2025) & Agent has strong tools and time budget & Latency,
non-determinism, repeated search loops & Prefer live search when local
files exist and freshness matters \\
Pre-indexing & Large repos need fast semantic or structural context &
Cursor and Sourcegraph/Amp claims (Cursor 2026; Sourcegraph 2025) &
Index is fresh, private enough, and amortized over many tasks &
Staleness, privacy, setup cost & Index when repo size and query
frequency justify maintenance \\
Repo maps & Summaries preserve structure under tight budgets & aider
tree-sitter repo map (Gauthier 2023) & Symbols/signatures are enough for
orientation & Missing details and runtime behavior & Use maps for
orientation, then expand to exact files \\
Skills and rules & Procedural knowledge should be progressively
disclosed & Agent Skills and CLAUDE.md practice (Anthropic 2025b;
HumanLayer (Kyle) 2025) & Knowledge is stable and reusable & Overbroad
instructions and hidden script risk & Use skills for procedures,
conventions, and reusable workflows \\
Always-on docs indexes & A compact project instruction file can reliably
point agents at versioned local docs & Vercel Next.js agent evals (Gao
2026a) & The index is small, accurate, and points to local docs the
agent can read & Context bloat; stale copied docs; agent ignores project
instruction files & Put high-value retrieval pointers in AGENTS.md, and
reserve skills for explicit workflows \\
Managed enterprise context & Org search needs ACLs, multi-repo scale,
and governance & Glean, Augment, Sourcegraph, DeepWiki-style systems &
Cloud or managed index is allowed & IP leakage, opaque ranking, vendor
lock-in & Route through managed systems only when policy permits \\
\end{longtable}

\subsection{Effect as the originating case}\label{sec-effect-origin}

Effect is the motivating case for ingraft and for this paper. Brown's
post makes a narrow but important source-vendoring claim: for an
idiom-heavy TypeScript framework, human-facing docs can be too lossy for
coding agents, while source, tests, examples, and module structure
preserve the patterns agents need (Brown 2026). The recommended route is
git subtree under a read-only reference directory, with agent
instructions and editor/tool ignores so humans are not overwhelmed by
vendored code.

The public article should not be treated as a controlled benchmark
report. It does not publish task definitions, model list, run counts, or
pass/fail tables. Its evidential role here is practitioner operating
experience: for Effect, the authors report that locally available source
is more useful to agents than fragmented web fetches, human docs, or
compiled \texttt{node\_modules} surfaces. That is enough to justify
source vendoring as a candidate route, not enough to make it the global
default.

This is exactly the router opening. If source vendoring wins for Effect,
the hidden variables are high package centrality, idiom density,
navigable TypeScript source, useful tests and examples, manageable
repository size, and need for version-specific source. The router
implication is not ``always vendor''; it is ``vendor when source
authority and idiom value outweigh repository impact and distractor
risk.''

\subsection{Always-on indexes versus invoked
skills}\label{sec-agents-md-evals}

Vercel's Next.js team published a useful negative result for the router
thesis: on a hardened Next.js 16 eval suite, an on-demand docs skill was
often not invoked and did not improve over the no-docs baseline, while
explicit instructions improved results but were wording-sensitive (Gao
2026a). Their best condition was a compressed \texttt{AGENTS.md} docs
index: the project file kept a small map of version-matched docs always
visible, and the agent read the actual docs from a local
\texttt{.next-docs/} directory when needed. The reported pass rates were
53\% for baseline, 53\% for default skill use, 79\% for skill plus
explicit instruction, and 100\% for the \texttt{AGENTS.md} docs index.

\figureVercelEval

This is vendor evidence, not an independent controlled study, and it is
specific to Next.js tasks and the harnesses tested. Still, it sharpens
the routing model: there are cases where the best route is not to load
full docs into context and not to rely on skill discovery, but to expose
a small always-on retrieval index that removes the agent's decision to
ask for help. Vercel's public \texttt{next-evals-oss} repository is also
a practical evaluation pattern: self-contained project fixtures with
\texttt{PROMPT.md}, withheld \texttt{EVAL.ts}, paired experiments with
and without \texttt{AGENTS.md}, and exported docs-impact deltas (Vercel
2026). The merged \texttt{@next/codemod\ agents-md} command shows the
implementation route: detect the framework version, fetch matching docs,
generate a compressed index, inject it between markers in
\texttt{AGENTS.md} or \texttt{CLAUDE.md}, and ignore the local docs
directory (Gao 2026b).

\subsection{Source is not always better than
docs}\label{sec-source-docs}

The Effect case is one operating zone, not a universal source-over-docs
rule. The source-vendoring argument is strongest when agents need
idioms, examples, type-level patterns, framework internals, and
implementation behavior. Docs are often the authority for public API
intent, migration guidance, deprecation warnings, and security
constraints.

\begin{longtable}[]{@{}
  >{\raggedright\arraybackslash}p{(\linewidth - 2\tabcolsep) * \real{0.5000}}
  >{\raggedright\arraybackslash}p{(\linewidth - 2\tabcolsep) * \real{0.5000}}@{}}
\caption{Source-versus-docs decision
rule.}\label{tbl-source-docs}\tabularnewline
\toprule\noalign{}
\begin{minipage}[b]{\linewidth}\raggedright
Task need
\end{minipage} & \begin{minipage}[b]{\linewidth}\raggedright
Prefer
\end{minipage} \\
\midrule\noalign{}
\endfirsthead
\toprule\noalign{}
\begin{minipage}[b]{\linewidth}\raggedright
Task need
\end{minipage} & \begin{minipage}[b]{\linewidth}\raggedright
Prefer
\end{minipage} \\
\midrule\noalign{}
\endhead
\bottomrule\noalign{}
\endlastfoot
Learn public API quickly & Docs and curated examples \\
Understand behavior edge case & Source and tests \\
Follow idioms & Source examples plus local usage \\
Avoid private/internal APIs & Official docs and public type surface \\
Debug package internals & Version-matched source \\
Migrate across versions & Changelog and migration guide, then local
usage \\
Verify generated code & Tests plus source \\
Understand package philosophy & Docs and design notes \\
\end{longtable}

\subsection{Agentic search versus indexing is a break-even
problem}\label{sec-break-even}

The unresolved question is not whether indexing or agentic search is
``right.'' It is when the cost of maintaining an index exceeds the cost
of repeated live search.

Agentic search wins when:

\[
C_\mathrm{toolCalls} + C_\mathrm{tokens} + C_\mathrm{latency}
<
C_\mathrm{indexBuild} + C_\mathrm{indexDrift} + C_\mathrm{privacy}
+ C_\mathrm{retrievalErrors}
\]

Indexing wins when repository size, query frequency, and cross-file
complexity are large enough to amortize index maintenance. This depends
on churn, local change frequency, language tooling, privacy rules, and
whether cross-repo search matters. The router should make that
break-even explicit rather than choosing a side globally.

\section{Router Design}\label{sec-router}

\subsection{Inputs and outputs}\label{sec-router-io}

A context router takes:

\begin{itemize}
\tightlist
\item
  task statement and task type;
\item
  repository state, including local changes and lockfiles;
\item
  knowledge object metadata;
\item
  model and tool capabilities;
\item
  privacy and trust policy;
\item
  history of failures or user overrides.
\end{itemize}

It outputs:

\begin{itemize}
\tightlist
\item
  selected route and alternatives considered;
\item
  context bundle or tool affordance;
\item
  provenance record;
\item
  trust class and intended use;
\item
  explanation and override path;
\item
  optional context-lock update.
\end{itemize}

\figureRouterArchitecture

\subsection{Strategy scoring}\label{sec-strategy-scoring}

The router can begin as a rule system, but the policy should expose
scoring features so users can understand why a route was chosen.

\begin{longtable}[]{@{}
  >{\raggedright\arraybackslash}p{(\linewidth - 10\tabcolsep) * \real{0.1667}}
  >{\raggedright\arraybackslash}p{(\linewidth - 10\tabcolsep) * \real{0.1667}}
  >{\raggedright\arraybackslash}p{(\linewidth - 10\tabcolsep) * \real{0.1667}}
  >{\raggedright\arraybackslash}p{(\linewidth - 10\tabcolsep) * \real{0.1667}}
  >{\raggedright\arraybackslash}p{(\linewidth - 10\tabcolsep) * \real{0.1667}}
  >{\raggedright\arraybackslash}p{(\linewidth - 10\tabcolsep) * \real{0.1667}}@{}}
\caption{Scoring features across route
families.}\label{tbl-scoring}\tabularnewline
\toprule\noalign{}
\begin{minipage}[b]{\linewidth}\raggedright
Feature
\end{minipage} & \begin{minipage}[b]{\linewidth}\raggedright
Vendoring
\end{minipage} & \begin{minipage}[b]{\linewidth}\raggedright
Lazy fetch
\end{minipage} & \begin{minipage}[b]{\linewidth}\raggedright
Docs
\end{minipage} & \begin{minipage}[b]{\linewidth}\raggedright
Index
\end{minipage} & \begin{minipage}[b]{\linewidth}\raggedright
Skill
\end{minipage} \\
\midrule\noalign{}
\endfirsthead
\toprule\noalign{}
\begin{minipage}[b]{\linewidth}\raggedright
Feature
\end{minipage} & \begin{minipage}[b]{\linewidth}\raggedright
Vendoring
\end{minipage} & \begin{minipage}[b]{\linewidth}\raggedright
Lazy fetch
\end{minipage} & \begin{minipage}[b]{\linewidth}\raggedright
Docs
\end{minipage} & \begin{minipage}[b]{\linewidth}\raggedright
Index
\end{minipage} & \begin{minipage}[b]{\linewidth}\raggedright
Skill
\end{minipage} \\
\midrule\noalign{}
\endhead
\bottomrule\noalign{}
\endlastfoot
Version fidelity & High if lockfile-resolved & High if resolved & Medium
to high & Varies & Low unless versioned \\
Token efficiency & Medium & Medium & High & High & High \\
Agent navigability & High & High & Medium & Medium/high & Low/medium \\
Setup cost & Medium/high & Low & Low & High & Low \\
Privacy & High & High & Varies & Varies & High \\
Freshness & Pinned & Pinned & Fresh but may mismatch & Drift risk &
Stable \\
Distractor risk & High if huge & Medium & Medium & Medium & Low \\
Best for & Deep idioms & Occasional inspection & Public API/migration &
Large repos & Conventions/procedures \\
\end{longtable}

Default policy rules follow from the table:

\begin{itemize}
\tightlist
\item
  Vendor when package centrality is high, source size is manageable,
  privacy matters, and idioms matter.
\item
  Lazy-fetch when a package is used occasionally but source-level
  inspection may be needed.
\item
  Docs-first when the task targets public API, migration, conceptual
  usage, security guidance, or newly released features.
\item
  Index when repository or package scale exceeds grep-friendly size, or
  cross-file relationships dominate.
\item
  Skill when knowledge is procedural, stable, reusable, and safe to
  package.
\item
  Agentic-search when local source is available and the model/tooling
  can inspect selectively.
\item
  Pack when the agent lacks filesystem/search tools or the task is
  one-shot and offline.
\end{itemize}

\subsection{Source authority and conflict
resolution}\label{sec-authority}

A router should ask not only ``where can I get information?'' but
``which source is authoritative for this question?''

\begin{longtable}[]{@{}
  >{\raggedright\arraybackslash}p{(\linewidth - 2\tabcolsep) * \real{0.5000}}
  >{\raggedright\arraybackslash}p{(\linewidth - 2\tabcolsep) * \real{0.5000}}@{}}
\caption{Source authority by question
type.}\label{tbl-authority}\tabularnewline
\toprule\noalign{}
\begin{minipage}[b]{\linewidth}\raggedright
Question
\end{minipage} & \begin{minipage}[b]{\linewidth}\raggedright
Best authority
\end{minipage} \\
\midrule\noalign{}
\endfirsthead
\toprule\noalign{}
\begin{minipage}[b]{\linewidth}\raggedright
Question
\end{minipage} & \begin{minipage}[b]{\linewidth}\raggedright
Best authority
\end{minipage} \\
\midrule\noalign{}
\endhead
\bottomrule\noalign{}
\endlastfoot
What version is installed? & Lockfile and package manager metadata \\
What public API is intended? & Official docs and public type surface \\
What behavior occurs? & Version-matched source and tests \\
What pattern does this repo prefer? & Local codebase and AGENTS.md \\
What should the agent never do? & AGENTS.md, security policy, human
instruction \\
What broke? & Failing test, log, trace, CI output \\
What changed between versions? & Changelog, release notes, migration
guide \\
What symbols reference this? & LSP, SCIP, source index, grep \\
\end{longtable}

Conflict resolution should preserve provenance. A useful agent-facing
message is not ``use zod docs''; it is ``this recommendation is based on
local package source resolved from \texttt{bun.lock}, not latest web
docs.'' That distinction prevents version skew and makes human override
possible.

\subsection{Provenance and context lockfiles}\label{sec-provenance}

The router should emit context plus provenance. Example bundle metadata
for source:

\begin{Shaded}
\begin{Highlighting}[]
\FunctionTok{context\_bundle}\KeywordTok{:}
\AttributeTok{  }\FunctionTok{source}\KeywordTok{:}\AttributeTok{ npm}
\AttributeTok{  }\FunctionTok{package}\KeywordTok{:}\AttributeTok{ effect}
\AttributeTok{  }\FunctionTok{version}\KeywordTok{:}\AttributeTok{ }\FloatTok{3.14.2}
\AttributeTok{  }\FunctionTok{route}\KeywordTok{:}\AttributeTok{ vendor}
\AttributeTok{  }\FunctionTok{resolved\_from}\KeywordTok{:}\AttributeTok{ bun.lock}
\AttributeTok{  }\FunctionTok{fetched\_at}\KeywordTok{:}\AttributeTok{ 2026{-}05{-}13}
\AttributeTok{  }\FunctionTok{trust\_class}\KeywordTok{:}\AttributeTok{ local{-}readonly{-}source}
\AttributeTok{  }\FunctionTok{mutable}\KeywordTok{:}\AttributeTok{ }\CharTok{false}
\AttributeTok{  }\FunctionTok{intended\_use}\KeywordTok{:}\AttributeTok{ reference{-}only}
\end{Highlighting}
\end{Shaded}

For docs:

\begin{Shaded}
\begin{Highlighting}[]
\FunctionTok{context\_bundle}\KeywordTok{:}
\AttributeTok{  }\FunctionTok{source}\KeywordTok{:}\AttributeTok{ context7}
\AttributeTok{  }\FunctionTok{package}\KeywordTok{:}\AttributeTok{ next}
\AttributeTok{  }\FunctionTok{version\_hint}\KeywordTok{:}\AttributeTok{ }\StringTok{"15.x"}
\AttributeTok{  }\FunctionTok{route}\KeywordTok{:}\AttributeTok{ docs}
\AttributeTok{  }\FunctionTok{fetched\_at}\KeywordTok{:}\AttributeTok{ 2026{-}05{-}13}
\AttributeTok{  }\FunctionTok{trust\_class}\KeywordTok{:}\AttributeTok{ remote{-}docs}
\AttributeTok{  }\FunctionTok{warning}\KeywordTok{:}\AttributeTok{ verify against lockfile and local source before editing}
\end{Highlighting}
\end{Shaded}

A \texttt{context.lock} file is the reproducibility bridge between
research and implementation. It records the strategy, source URL,
package version, commit, docs version, index version, skill version,
trust class, generated summaries, timestamps, and override reasons.

\begin{Shaded}
\begin{Highlighting}[]
\KeywordTok{[packages.effect]}
\DataTypeTok{strategy} \OperatorTok{=} \StringTok{"vendor"}
\DataTypeTok{version} \OperatorTok{=} \StringTok{"3.14.2"}
\DataTypeTok{source} \OperatorTok{=} \StringTok{"github:Effect{-}TS/effect"}
\DataTypeTok{commit} \OperatorTok{=} \StringTok{"abc123"}
\DataTypeTok{reason} \OperatorTok{=} \StringTok{"core{-}framework allowlist; imported in 42 files"}
\DataTypeTok{trust} \OperatorTok{=} \StringTok{"local{-}readonly"}

\KeywordTok{[packages.zod]}
\DataTypeTok{strategy} \OperatorTok{=} \StringTok{"docs{-}first"}
\DataTypeTok{version} \OperatorTok{=} \StringTok{"4.1.0"}
\DataTypeTok{docs} \OperatorTok{=} \StringTok{"context7:zod@4.1"}
\DataTypeTok{fallback} \OperatorTok{=} \StringTok{"fetch{-}source"}
\DataTypeTok{reason} \OperatorTok{=} \StringTok{"public API task; source size low; docs available"}
\DataTypeTok{trust} \OperatorTok{=} \StringTok{"remote{-}docs"}

\KeywordTok{[tools.mgrep]}
\DataTypeTok{enabled} \OperatorTok{=} \ConstantTok{false}
\DataTypeTok{reason} \OperatorTok{=} \StringTok{"privacy=strict{-}local"}
\end{Highlighting}
\end{Shaded}

\begin{longtable}[]{@{}ll@{}}
\caption{Core fields for a context
lockfile.}\label{tbl-lockfields}\tabularnewline
\toprule\noalign{}
Lockfile field & Purpose \\
\midrule\noalign{}
\endfirsthead
\toprule\noalign{}
Lockfile field & Purpose \\
\midrule\noalign{}
\endhead
\bottomrule\noalign{}
\endlastfoot
\texttt{object\_id} & Stable identity for the routed artifact \\
\texttt{kind} & Context-object class \\
\texttt{source\_uri} & Upstream origin \\
\texttt{revision} & Commit, version, timestamp, or content hash \\
\texttt{authority} & Canonical, upstream, derived, advisory, or
unverified \\
\texttt{materialization} & Committed, cache, mount, live, or pack \\
\texttt{representation} & Raw, summary, metadata, index, or executable
tool \\
\texttt{privacy\_class} & Repo-safe, local-only, restricted, or
remote \\
\texttt{integrity} & Hash, signature, or provenance attestation \\
\texttt{phase\_policy} & Preferred routes by task phase \\
\texttt{refresh\_policy} & Pinned, manual refresh, TTL, or live \\
\texttt{transform\_chain} & Export, normalization, and summarization
steps \\
\end{longtable}

Provenance supports version fidelity, audit logs, privacy review,
benchmark reproducibility, and human overrides.

\subsection{Trust policy}\label{sec-security}

Security is central to routing, not an appendix. Every route has a
different attack surface.

\begin{longtable}[]{@{}
  >{\raggedright\arraybackslash}p{(\linewidth - 2\tabcolsep) * \real{0.5000}}
  >{\raggedright\arraybackslash}p{(\linewidth - 2\tabcolsep) * \real{0.5000}}@{}}
\caption{Route-specific security
risks.}\label{tbl-security}\tabularnewline
\toprule\noalign{}
\begin{minipage}[b]{\linewidth}\raggedright
Route
\end{minipage} & \begin{minipage}[b]{\linewidth}\raggedright
Risk
\end{minipage} \\
\midrule\noalign{}
\endfirsthead
\toprule\noalign{}
\begin{minipage}[b]{\linewidth}\raggedright
Route
\end{minipage} & \begin{minipage}[b]{\linewidth}\raggedright
Risk
\end{minipage} \\
\midrule\noalign{}
\endhead
\bottomrule\noalign{}
\endlastfoot
Docs retrieval & Prompt injection, outdated examples, wrong version \\
MCP server & Tool overpermission, data exfiltration, remote execution \\
Skills & Executable scripts, hidden instructions, supply-chain
attacks \\
Vendored source & Malicious comments, enormous distractor surface \\
Repo pack & Secret exposure and uncontrolled redistribution \\
Cloud index & IP leakage and opaque retention \\
Agentic search & Command misuse and accidental broad reads \\
Generated summaries & Untraceable hallucinated compression \\
\end{longtable}

A practical router should support policy modes such as:

\begin{itemize}
\tightlist
\item
  \texttt{strict-local}: no remote retrieval or cloud index;
\item
  \texttt{local-executable-denied}: local files allowed, bundled scripts
  blocked;
\item
  \texttt{remote-docs-allowed}: remote docs allowed but quoted as
  untrusted;
\item
  \texttt{remote-code-denied}: no remote source fetch;
\item
  \texttt{cloud-index-allowed}: managed indexes permitted;
\item
  \texttt{scripts-require-confirmation}: skills and MCP tools cannot
  execute without confirmation;
\item
  \texttt{untrusted-context-quoted}: external context is evidence, not
  instruction.
\end{itemize}

The ethical case for an open router follows from these requirements:
transparent routing decisions, user override, local-first defaults,
provenance logs, auditable lockfiles, no silent cloud upload,
standards-based portability, and avoidance of vendor lock-in.

\subsection{Explainability requirements}\label{sec-explainability}

A useful routing explanation includes:

\begin{itemize}
\tightlist
\item
  selected strategy;
\item
  alternatives considered;
\item
  reason selected;
\item
  risks;
\item
  override command or config;
\item
  provenance;
\item
  expected cost.
\end{itemize}

Example:

\begin{quote}
Routed \texttt{effect} to \texttt{vendor} because it is imported in 47
files, appears in the core-framework allowlist, source size is below 25
MB, and privacy policy is \texttt{strict-local}. Alternatives:
docs-first rejected because the task requires idiomatic patterns; remote
index rejected by privacy policy. Override with
\texttt{ingraft\ route\ effect\ -\/-strategy\ fetch}.
\end{quote}

Explanations are not decoration. They are how users correct the policy.

\subsection{Worked routing examples}\label{sec-examples}

\textbf{Daily framework dependency.} A project uses Effect in 80 files.
The task is to implement a new service using local idioms. The router
vendors Effect source read-only, emits AGENTS.md guidance, and tells the
agent to inspect local examples before docs.

\textbf{Rare utility package.} The project uses \texttt{nanoid} in one
file and the task is unrelated. The router does nothing. If the agent
asks, it fetches source lazily.

\textbf{API migration.} A project upgrades Next.js from 14 to 15. The
router prioritizes version-specific migration docs and changelog, then
local usage search, then source fallback.

\textbf{Large monorepo.} The task touches one service in a 500K-file
monorepo. The router avoids full packing, queries structural code search
or SCIP/LSP, and passes only relevant files.

\textbf{Privacy-restricted enterprise repo.} Cloud retrieval is
disabled. The router uses local source, local docs, local symbolic
search, and blocks remote MCPs.

\section{Standards and Interoperability}\label{sec-standards}

Standards matter because routing should not be trapped inside one agent
vendor. The standards landscape is increasingly centered around AAIF
projects such as MCP, goose, and AGENTS.md, alongside adjacent standards
and conventions such as Agent Skills, ACP, A2A, llms.txt, and SCIP. It
would be wrong to imply that all of these have one governance story.

The Linux Foundation announced the Agentic AI Foundation on December 9,
2025, with founding project contributions including Anthropic's Model
Context Protocol, Block's goose, and OpenAI's AGENTS.md (Linux
Foundation 2025). The same announcement describes MCP as a protocol for
connecting models to tools, data, and applications, and AGENTS.md as a
project-guidance standard for coding agents. Agent Skills has its own
specification and progressive-disclosure model (Agent Skills Working
Group 2025; Anthropic 2025b). ACP, A2A, llms.txt, and SCIP have separate
origins and adoption paths (Zed Industries 2025; Google and Linux
Foundation 2025; Howard 2024; Sourcegraph 2024).

\begin{longtable}[]{@{}
  >{\raggedright\arraybackslash}p{(\linewidth - 2\tabcolsep) * \real{0.5000}}
  >{\raggedright\arraybackslash}p{(\linewidth - 2\tabcolsep) * \real{0.5000}}@{}}
\caption{Standards and their routing
role.}\label{tbl-standards-role}\tabularnewline
\toprule\noalign{}
\begin{minipage}[b]{\linewidth}\raggedright
Standard or convention
\end{minipage} & \begin{minipage}[b]{\linewidth}\raggedright
Router role
\end{minipage} \\
\midrule\noalign{}
\endfirsthead
\toprule\noalign{}
\begin{minipage}[b]{\linewidth}\raggedright
Standard or convention
\end{minipage} & \begin{minipage}[b]{\linewidth}\raggedright
Router role
\end{minipage} \\
\midrule\noalign{}
\endhead
\bottomrule\noalign{}
\endlastfoot
AGENTS.md & Declares project-level routing policy and agent
instructions \\
MCP & Exposes router tools to multiple agents \\
Agent Skills & Packages strategy-specific procedures with progressive
disclosure \\
ACP & Lets editors invoke agents that call the router \\
A2A & Lets agents coordinate or delegate context work \\
SCIP/LSP & Supplies graph and symbol context \\
llms.txt & Provides docs-discovery hints, low authority unless
verified \\
\end{longtable}

The portable stack is therefore not ``one standard to rule them all.''
It is: AGENTS.md for durable project policy, MCP for tool access, Agent
Skills for procedural knowledge, and language/index standards for
structural retrieval.

\section{Evaluation Agenda}\label{sec-eval}

The router thesis needs a router eval. The key research question is:

\begin{quote}
Does a task-aware context router outperform fixed context strategies
under equal model, budget, and tool conditions?
\end{quote}

\subsection{Strategies to compare}\label{sec-eval-strategies}

\begin{longtable}[]{@{}
  >{\raggedright\arraybackslash}p{(\linewidth - 2\tabcolsep) * \real{0.5000}}
  >{\raggedright\arraybackslash}p{(\linewidth - 2\tabcolsep) * \real{0.5000}}@{}}
\caption{Fixed strategies and router
condition.}\label{tbl-eval-strategies}\tabularnewline
\toprule\noalign{}
\begin{minipage}[b]{\linewidth}\raggedright
Strategy
\end{minipage} & \begin{minipage}[b]{\linewidth}\raggedright
Description
\end{minipage} \\
\midrule\noalign{}
\endfirsthead
\toprule\noalign{}
\begin{minipage}[b]{\linewidth}\raggedright
Strategy
\end{minipage} & \begin{minipage}[b]{\linewidth}\raggedright
Description
\end{minipage} \\
\midrule\noalign{}
\endhead
\bottomrule\noalign{}
\endlastfoot
No extra context & Baseline agent with repo only \\
Docs-only & Version-specific docs retrieval where available \\
Always-on docs index & AGENTS.md/CLAUDE.md points to local versioned
docs \\
Source vendoring & Dependency source in workspace \\
Lazy source fetch & Package source fetched on demand \\
Repo pack & Repomix-style snapshot \\
Semantic index & Embedding or code-search retrieval \\
Structural index & AST, LSP, SCIP, call/reference graph retrieval \\
Agentic search & Bash/grep/read loop \\
Skills/rules & Procedural guidance only \\
Router & Chooses among the above \\
\end{longtable}

\figureEvalLattice

\subsection{Task families}\label{sec-eval-tasks}

\begin{longtable}[]{@{}
  >{\raggedright\arraybackslash}p{(\linewidth - 2\tabcolsep) * \real{0.5000}}
  >{\raggedright\arraybackslash}p{(\linewidth - 2\tabcolsep) * \real{0.5000}}@{}}
\caption{Task families for router
evaluation.}\label{tbl-eval-tasks}\tabularnewline
\toprule\noalign{}
\begin{minipage}[b]{\linewidth}\raggedright
Task family
\end{minipage} & \begin{minipage}[b]{\linewidth}\raggedright
Example
\end{minipage} \\
\midrule\noalign{}
\endfirsthead
\toprule\noalign{}
\begin{minipage}[b]{\linewidth}\raggedright
Task family
\end{minipage} & \begin{minipage}[b]{\linewidth}\raggedright
Example
\end{minipage} \\
\midrule\noalign{}
\endhead
\bottomrule\noalign{}
\endlastfoot
Unknown library use & Integrate unfamiliar package \\
Framework idiom task & Add route, middleware, schema, service, or
hook \\
API migration & Upgrade package version \\
Bugfix with stack trace & Failure crosses dependency boundary \\
Multi-file feature & Edit across local repo \\
Security-sensitive change & Avoid unsafe API or pattern \\
Performance/debugging task & Inspect implementation behavior \\
Test repair & Fix tests requiring project conventions \\
\end{longtable}

Stratification variables should include package popularity, source size,
docs quality, framework versus utility package, typed versus untyped
ecosystem, repo size, public API versus internals, package version
freshness, model class, and privacy setting.

\subsection{Metrics, ablations, and
falsification}\label{sec-eval-metrics}

\begin{longtable}[]{@{}ll@{}}
\caption{Router evaluation
metrics.}\label{tbl-eval-metrics}\tabularnewline
\toprule\noalign{}
Metric & Why it matters \\
\midrule\noalign{}
\endfirsthead
\toprule\noalign{}
Metric & Why it matters \\
\midrule\noalign{}
\endhead
\bottomrule\noalign{}
\endlastfoot
pass@1 / pass@k & Task success \\
Tests passed & Functional correctness \\
Human review score & Maintainability and style \\
Tokens used & Context efficiency \\
Wall-clock latency & Practical usability \\
Tool calls & Loop complexity \\
Retrieval precision & Context quality \\
Provenance coverage & Auditability \\
Version match rate & Source/doc correctness \\
Regression rate & Safety \\
Cost per solved task & Operational relevance \\
\end{longtable}

Ablations should remove model awareness, privacy constraints, vendoring,
docs, graph/AST retrieval, stale-context detection, and
failure-triggered strategy switching. Stress tests should include stale
docs, distractor context, wrong package versions, and small versus
frontier models.

The thesis should be considered weakened if:

\begin{itemize}
\tightlist
\item
  one fixed strategy beats the router across most task/model/repo
  settings;
\item
  router explanations do not help users override bad choices;
\item
  routing overhead cancels out quality gains;
\item
  context strategy matters less than model choice once modern agents are
  used;
\item
  simpler agentic search dominates retrieval, indexing, and vendoring
  under realistic budgets.
\end{itemize}

Strong papers say what would falsify them. These are the falsification
tests.

\section{Case Study: ingraft}\label{sec-ingraft}

This paper is motivated by ingraft, an open-source CLI for package and
repository context routing (Brunner 2026). The case study should be read
as an implementation path, not as empirical proof that the router wins.

\subsection{Scope}\label{sec-ingraft-scope}

ingraft's natural first slice is package knowledge:

\begin{itemize}
\tightlist
\item
  inspect package manager state and lockfiles;
\item
  route core dependencies to read-only vendored source when justified;
\item
  route long-tail dependencies to lazy fetch;
\item
  route public API or migration tasks to docs-first;
\item
  expose provenance and route explanations;
\item
  write AGENTS.md guidance so agents treat vendored source as reference
  material and do not import from it.
\end{itemize}

This maps to the project rule already used in this repository: vendored
repositories under \texttt{vendor/} are read-only reference material,
not application code.

\subsection{From package routing to knowledge
routing}\label{sec-knowledge-routing}

The same design generalizes to non-code knowledge, but with sharper
privacy and authority rules. A Linear ticket, Notion decision page,
OpenAPI spec, CI log, or Figma design note can be treated as a context
object. The router can choose to vendor it as committed text, fetch it
into a local cache, mount it read-only, query it through an MCP server,
or summarize it into a skill.

The important design point is not ``everything belongs in
\texttt{vendor/}.'' It is that every artifact needs an explicit route,
provenance record, privacy class, and authority label.

Example workspace layout:

\begin{Shaded}
\begin{Highlighting}[]
\NormalTok{vendor/                         \# committed, read{-}only reference material}
\NormalTok{  github/effect{-}ts/effect/       \# source snapshot}
\NormalTok{  notion/team{-}handbook/index.md  \# exported page}
\NormalTok{  linear/DEV{-}42/index.md         \# ticket snapshot}
\NormalTok{  openapi/payments{-}api.json      \# pinned API spec}

\NormalTok{.ingraft/}
\NormalTok{  config.toml                    \# committed routing policy}
\NormalTok{  context.lock                   \# committed or reviewed provenance lock}
\NormalTok{  cache/                         \# gitignored lazy fetches}
\NormalTok{  mounts/                        \# gitignored live mounts}
\end{Highlighting}
\end{Shaded}

\subsection{ingraft limitations}\label{sec-ingraft-limitations}

The current evidence supports ingraft as a promising implementation of a
repo-resident router, not as an established empirical winner. The next
step is the router eval in Section~\ref{sec-eval}. ingraft should also
treat security as a first class feature: signed connector metadata,
local-first defaults, privacy-class checks, provenance logs, and
explicit user confirmation before executable skills or connector scripts
run.

\section{Threats to Validity and Counterexamples}\label{sec-threats}

\textbf{Benchmark validity.} SWE-bench, LongCodeBench, CodeRAG-Bench,
RepoEval, CrossCodeEval, and vendor evals measure different task
distributions. Improvements on one may not transfer to another.

\textbf{Model drift.} Results from Claude 3.5, GPT-4o, Qwen2.5, Gemini
1.5, and other models may not extrapolate to newer systems. A router
should use model-aware calibration, not hardcoded thresholds.

\textbf{Tooling confounds.} Agent performance reflects model, scaffold,
tools, prompts, editing loop, test harness, retry policy, and human
oversight. It is hard to isolate context strategy alone.

\textbf{Vendor telemetry.} Cursor, Augment, Cognition, Greptile, and
similar numbers may be directionally useful but should not be treated as
independent evidence unless methodology is public.

\textbf{Publication bias.} Practitioner essays report successful
workflows more often than failed context strategies.

\textbf{Temporal instability.} Tool features, standards adoption,
pricing, security posture, and model context behavior change quickly.
Dated snapshots should be clearly labeled.

Routing may not help much in very small repos where the whole context
fits, tasks dominated by model reasoning rather than repo knowledge,
highly standardized frameworks where docs are enough, environments where
the agent already has excellent local tool use, security-constrained
settings where remote retrieval is forbidden and vendoring is
impossible, product-judgment tasks, repos with poor tests, languages
with weak structural tooling, generated or minified codebases, and
monorepos whose package boundaries do not match architectural
boundaries.

\section{Predictions}\label{sec-predictions}

The router thesis makes testable predictions:

\begin{enumerate}
\def\labelenumi{\arabic{enumi}.}
\tightlist
\item
  Hybrid systems will outperform single-context systems under diverse
  task distributions.
\item
  The optimal route will vary by package centrality, not merely repo
  size.
\item
  Source vendoring will help most on idiom-heavy frameworks and least on
  stable utility libraries.
\item
  Docs-first retrieval will outperform source-first on migration tasks
  and public API usage.
\item
  Graph/AST retrieval will show larger gains as cross-file dependency
  depth increases.
\item
  Aggressive retrieval will hurt more often on distractor-heavy tasks
  than on tightly scoped tasks.
\item
  User-visible routing explanations will reduce repeated failure loops
  by making overrides possible.
\end{enumerate}

\section{Conclusion}\label{sec-conclusion}

The 2026 context-engineering landscape is not a contest between ``vendor
source'' and ``use RAG.'' It is a routing problem. Context windows
remain attention scarce, retrieval is conditional, structure matters,
authority matters, freshness matters, and privacy often constrains the
route before quality optimization begins.

The strongest form of the thesis is:

\begin{quote}
Context engineering for coding agents is context allocation under
attention, authority, freshness, privacy, trust, and cost constraints.
Existing tools each solve one region of the design space. A router is
the abstraction that makes those partial solutions composable.
\end{quote}

An open, repo-resident routing layer is a promising default because it
is inspectable, overrideable, local-first, and portable across agent
vendors. Demonstrating its empirical advantage over fixed strategies
remains the next evaluation task.

\begin{tcolorbox}[enhanced jigsaw, arc=.35mm, bottomrule=.15mm, bottomtitle=1mm, breakable, colback=white, colbacktitle=paper-callout-tip-color!10!white, colframe=paper-callout-tip-color-frame, coltitle=black, left=2mm, leftrule=.75mm, opacityback=0, opacitybacktitle=0.6, rightrule=.15mm, title=\textcolor{paper-callout-tip-color}{\faLightbulb}\hspace{0.5em}{Citation}, titlerule=0mm, toprule=.15mm, toptitle=1mm]

If you reference this work in academic or industry writing:

\begin{quote}
Brunner, G. \emph{Context Routing for AI Coding Agents: A Research
Agenda for Context Allocation Under Attention, Authority, Freshness,
Privacy, and Cost Constraints.} ingraft project, May 2026.
\url{https://ingraft.dev/whitepaper}
\end{quote}

Issues, corrections, and disagreements are welcome at
\href{https://github.com/gunta/ingraft/issues}{github.com/gunta/ingraft/issues}.

\end{tcolorbox}

\emph{This whitepaper is released under
\href{https://creativecommons.org/licenses/by/4.0/}{CC BY 4.0}. Quote
freely with attribution.}

\section{Appendix A: Representative Tool Snapshot, May
2026}\label{sec-tools}

This appendix preserves the tool-catalog value without letting the
catalog drive the main argument. It is a representative snapshot, not a
complete inventory.

\figureToolLandscapeCounts

\begin{longtable}[]{@{}
  >{\raggedright\arraybackslash}p{(\linewidth - 4\tabcolsep) * \real{0.3333}}
  >{\raggedright\arraybackslash}p{(\linewidth - 4\tabcolsep) * \real{0.3333}}
  >{\raggedright\arraybackslash}p{(\linewidth - 4\tabcolsep) * \real{0.3333}}@{}}
\caption{Representative May 2026 tool
snapshot.}\label{tbl-tool-snapshot}\tabularnewline
\toprule\noalign{}
\begin{minipage}[b]{\linewidth}\raggedright
Category
\end{minipage} & \begin{minipage}[b]{\linewidth}\raggedright
Representative tools
\end{minipage} & \begin{minipage}[b]{\linewidth}\raggedright
Routing lesson
\end{minipage} \\
\midrule\noalign{}
\endfirsthead
\toprule\noalign{}
\begin{minipage}[b]{\linewidth}\raggedright
Category
\end{minipage} & \begin{minipage}[b]{\linewidth}\raggedright
Representative tools
\end{minipage} & \begin{minipage}[b]{\linewidth}\raggedright
Routing lesson
\end{minipage} \\
\midrule\noalign{}
\endhead
\bottomrule\noalign{}
\endlastfoot
IDE-native context engines & Cursor, Windsurf, GitHub Copilot, Augment,
JetBrains AI/Junie, Cline, Zed, Roo Code, Continue, opencode, aider,
Sourcegraph Amp, Amazon Q, Devin, OpenHands, Factory Droid, Kiro, Google
Jules, Warp Agent Mode, Codebuff, Plandex & IDEs mix index, manual
mention, agentic search, and rules; the router should expose which layer
is active \\
Standalone code search/indexing & Sourcebot, Zoekt, Livegrep, Hound,
ast-grep, Comby, Bloop, Greptile, Glean Code Search, Unblocked, Tabby,
Serena MCP, CocoIndex, Indexify & Search wins when scale or cross-file
structure overwhelms live grep \\
Repo packaging/snapshots & Repomix, gitingest, code2prompt,
files-to-prompt, yek, ai-digest, gpt-repository-loader & Packs are
useful one-shot snapshots but weak on freshness and position control \\
MCP context servers & Context7, filesystem MCP, GitHub MCP, GitLab MCP,
DeepWiki MCP, fff, mgrep, Exa MCP, Tavily MCP, Perplexity MCP,
MCPfinder, Glama, Smithery, mcp.so, Memory Bank, Mem0, Letta, Mnemos,
Pieces LTM-2, OpenSrc & MCP turns context into tools; tool permissions
and provenance become router concerns \\
Embeddings/RAG components & voyage-code-3, Jina code embeddings, Jina v2
base code, mxbai-embed-large, CodeBERT, GraphCodeBERT, UniXcoder,
CodeT5+, StarCoder, CodeSage, ChromaDB, Qdrant, Pinecone, Weaviate,
Milvus, pgvector & Embeddings are substrates, not policies; quality
depends on selector, representation, and task \\
Multi-repo/org context & RepoSwarm, Backstage TechDocs, Roadie RAG,
Glean, Sourcegraph, DeepWiki, Unblocked & Org context requires ACLs,
provenance, and privacy-first routing \\
Adjacent safety/workflow & Qodo Merge, GitGuardian ggshield AI Hook,
MCP-omnisearch, GPTScript, Jujutsu, agentjj, jj-skill & Context routing
touches review, secrets, web search, shared prompts, and VCS
workflows \\
\end{longtable}

\section{Appendix B: Paper Matrix}\label{sec-paper-matrix}

\begin{longtable}[]{@{}
  >{\raggedright\arraybackslash}p{(\linewidth - 10\tabcolsep) * \real{0.1667}}
  >{\raggedright\arraybackslash}p{(\linewidth - 10\tabcolsep) * \real{0.1667}}
  >{\raggedright\arraybackslash}p{(\linewidth - 10\tabcolsep) * \real{0.1667}}
  >{\raggedright\arraybackslash}p{(\linewidth - 10\tabcolsep) * \real{0.1667}}
  >{\raggedright\arraybackslash}p{(\linewidth - 10\tabcolsep) * \real{0.1667}}
  >{\raggedright\arraybackslash}p{(\linewidth - 10\tabcolsep) * \real{0.1667}}@{}}
\caption{Academic and technical works mapped to router
implications.}\label{tbl-paper-matrix}\tabularnewline
\toprule\noalign{}
\begin{minipage}[b]{\linewidth}\raggedright
Work
\end{minipage} & \begin{minipage}[b]{\linewidth}\raggedright
Task
\end{minipage} & \begin{minipage}[b]{\linewidth}\raggedright
Context type
\end{minipage} & \begin{minipage}[b]{\linewidth}\raggedright
Main finding
\end{minipage} & \begin{minipage}[b]{\linewidth}\raggedright
Router implication
\end{minipage} & \begin{minipage}[b]{\linewidth}\raggedright
Evidence strength
\end{minipage} \\
\midrule\noalign{}
\endfirsthead
\toprule\noalign{}
\begin{minipage}[b]{\linewidth}\raggedright
Work
\end{minipage} & \begin{minipage}[b]{\linewidth}\raggedright
Task
\end{minipage} & \begin{minipage}[b]{\linewidth}\raggedright
Context type
\end{minipage} & \begin{minipage}[b]{\linewidth}\raggedright
Main finding
\end{minipage} & \begin{minipage}[b]{\linewidth}\raggedright
Router implication
\end{minipage} & \begin{minipage}[b]{\linewidth}\raggedright
Evidence strength
\end{minipage} \\
\midrule\noalign{}
\endhead
\bottomrule\noalign{}
\endlastfoot
N. F. Liu et al. (2024) & QA & Long context & Position matters &
Prioritize placement and bundle size & Peer-reviewed journal \\
{Hsieh et al.} (2024) & Synthetic long-context & Multi-hop long context
& Effective length varies by model/task & Calibrate per model &
Peer-reviewed conference \\
Rando et al. (2025) & Code QA and repair & Long repo context & Long
context degrades unevenly on real code tasks & Avoid stuffing whole
repos; evaluate by task & Preprint/benchmark \\
Hong et al. (2025) & Long-context reliability & Distractors, semantic
needles, LongMemEval, repeated words & Increasing input length alone can
make outputs less reliable & Treat context rot as a first-class routing
cost & Technical report \\
Modarressi et al. (2025) & Long-context retrieval & Non-literal
needles & Literal matching overstates long-context capability & Prefer
semantic retrieval tests, not only lexical NIAH & Conference/preprint \\
Fu et al. (2025) & Missing-information detection & Omitted context &
Models struggle to identify absences even when retrieval looks easy &
Do not equate retrieval with comprehension & Conference/preprint \\
{Z. Z. Wang et al.} (2025) & Code generation & Retrieved docs/code &
Retrieval helps unevenly and retrievers miss useful contexts & Adaptive
retrieval, not fixed top-k & Peer-reviewed findings \\
Thakur et al. (2021) & Information retrieval & Heterogeneous retrieval
tasks & No retriever dominates zero-shot across domains & Route by task
distribution and retriever failure mode & Peer-reviewed benchmark \\
Khattab and Zaharia (2020) & Passage search & Late interaction &
Token-level matching improves retrieval while allowing precomputation &
Use late interaction when dense vectors are too lossy & Peer-reviewed
conference \\
Formal et al. (2021) & Information retrieval & Learned sparse
retrieval & Sparse expansion bridges lexical auditability and semantic
matching & Include sparse search in hybrid routing & Peer-reviewed
conference \\
Husain et al. (2019) & Code search & Natural-language to code & Code
search has a language gap distinct from generic text search & Evaluate
code retrieval separately from text retrieval & Benchmark/preprint \\
Li et al. (2024) & Code information retrieval & Multi-task code IR &
Code retrieval decomposes into different task families & Do not tune one
retriever for all code questions & Preprint/benchmark \\
Jiang et al. (2025) & Issue localization & Repository graph search &
Graph-guided search can outperform flatter localization routes & Add
graph search as an agentic fallback & Preprint \\
{``{U-NIAH}''} (2025) & Long-context/RAG & Synthetic needles/distractors
& RAG helps some models and hurts others & Model-aware routing &
Journal/preprint \\
F. Zhang et al. (2023) & Code completion & Iterative repo retrieval &
Repository context improves completion & Expand context iteratively &
Peer-reviewed conference \\
{Ding et al.} (2023) & Code completion & Cross-file context &
Single-file context is insufficient & Use repo-aware retrieval &
Peer-reviewed benchmark \\
{Liang et al.} (2024) & Code completion & Fused dual context & Context
quality and latency trade off & Rank, truncate, and score context &
Preprint \\
{Phan et al.} (2025) & Code completion & Semantic graph & Graph
retrieval outperforms flat baselines & Prefer structural retrieval when
possible & Workshop/preprint \\
{Y. (Sherry). Zhang et al.} (2025) & Code RAG & AST chunks & AST
chunking improves over line chunks & Structural chunking & Preprint \\
Brown (2026) & Practitioner source-vendoring report & Effect source tree
& Locally vendored source gives agents implementation structure, tests,
and examples that docs and web fetches omit & Treat source vendoring as
a candidate route for high-idiom, high-centrality packages &
Practitioner/vendor evidence \\
Gao (2026a) & Framework-agent evals & AGENTS.md docs index versus skills
& Always-on compressed docs index outperformed on-demand skill use in
their Next.js suite & Evaluate selector/timing, not only content quality
& Vendor engineering eval \\
Shi et al. (2023) & Reasoning & Irrelevant context & Distractors reduce
performance & Penalize distractor risk & Peer-reviewed conference \\
\end{longtable}

\section*{References}\label{references}
\addcontentsline{toc}{section}{References}

\protect\phantomsection\label{refs}
\begin{CSLReferences}{1}{1}
\bibitem[\citeproctext]{ref-agentskills2025spec}
Agent Skills Working Group. 2025. {``Agent Skills Specification.''}
\url{https://agentskills.io/specification}.

\bibitem[\citeproctext]{ref-anthropic2025contexteng}
Anthropic. 2025a. {``Effective Context Engineering for {AI} Agents.''}
September 29.
\url{https://www.anthropic.com/engineering/effective-context-engineering-for-ai-agents}.

\bibitem[\citeproctext]{ref-anthropic2025skills}
Anthropic. 2025b. {``Equipping Agents for the Real World with Agent
Skills.''} October 16.
\url{https://www.anthropic.com/engineering/equipping-agents-for-the-real-world-with-agent-skills}.

\bibitem[\citeproctext]{ref-asai2024selfrag}
Asai, Akari, Zeqiu Wu, Yizhong Wang, Avirup Sil, and Hannaneh
Hajishirzi. 2024. {``Self-{RAG}: Learning to Retrieve, Generate, and
Critique Through Self-Reflection.''} \emph{ICLR 2024}.
\url{https://arxiv.org/abs/2310.11511}.

\bibitem[\citeproctext]{ref-cline2025nondex}
Baumann, Nick. 2025. {``Why {Cline} Doesn't Index Your Codebase (and Why
That's a Good Thing).''} May 27.
\url{https://cline.bot/blog/why-cline-doesnt-index-your-codebase-and-why-thats-a-good-thing}.

\bibitem[\citeproctext]{ref-brown2026effect}
Brown, Maxwell. 2026. {``The One Weird Git Trick That Makes Coding
Agents More {Effect}-Ive.''} May 11.
\url{https://effect.website/blog/the-one-weird-git-trick-that-makes-coding-agents-more-effect-ive/}.

\bibitem[\citeproctext]{ref-ingraft2026}
Brunner, Gunther. 2026. {``{ingraft}: The Context Router for {AI} Coding
Agents.''} \url{https://ingraft.dev/}.

\bibitem[\citeproctext]{ref-cursor2026indexing}
Cursor. 2026. {``Securely Indexing Large Codebases.''} January 27.
\url{https://cursor.com/blog/secure-codebase-indexing}.

\bibitem[\citeproctext]{ref-ding2023crosscodeeval}
{Ding, Yangruibo et al.} 2023. {``{CrossCodeEval}: A Diverse and
Multilingual Benchmark for Cross-File Code Completion.''} \emph{NeurIPS
2023 Datasets and Benchmarks Track}.
\url{https://arxiv.org/abs/2310.11248}.

\bibitem[\citeproctext]{ref-ding2024cocomic}
Ding, Yangruibo, Zijian Wang, Wasi Uddin Ahmad, et al. 2024.
{``{CoCoMIC}: Code Completion by Jointly Modeling in-File and Cross-File
Context.''} \emph{Proceedings of LREC-COLING 2024}.
\url{https://arxiv.org/abs/2212.10007}.

\bibitem[\citeproctext]{ref-vercel2026agentsmd}
Gao, Jude. 2026a. {``{AGENTS.md} Outperforms Skills in Our Agent
Evals.''} January 27.
\url{https://vercel.com/blog/agents-md-outperforms-skills-in-our-agent-evals}.

\bibitem[\citeproctext]{ref-vercel2026nextagentspr}
Gao, Jude. 2026b. {``Feat(next-Codemod): Add Agents-Md Command for {AI}
Coding Agents.''} January 24.
\url{https://github.com/vercel/next.js/pull/88961}.

\bibitem[\citeproctext]{ref-gauthier2023repomap}
Gauthier, Paul. 2023. {``Building a Better Repository Map with
Tree-Sitter.''} October 22.
\url{https://aider.chat/2023/10/22/repomap.html}.

\bibitem[\citeproctext]{ref-a2a2025spec}
Google, and Linux Foundation. 2025. {``Agent-to-Agent Protocol
({A2A}).''} \url{https://a2a-protocol.org/latest/}.

\bibitem[\citeproctext]{ref-llmstxt2024spec}
Howard, Jeremy. 2024. {``The {llms.txt} Specification.''} September 3.
\url{https://llmstxt.org/}.

\bibitem[\citeproctext]{ref-hsieh2024ruler}
{Hsieh, Cheng-Ping et al.} 2024. {``{RULER}: What's the Real Context
Size of Your Long-Context Language Models?''} \emph{COLM 2024}.
\url{https://arxiv.org/abs/2404.06654}.

\bibitem[\citeproctext]{ref-humanlayer2025claudemd}
HumanLayer (Kyle). 2025. {``Writing a Good {CLAUDE.md}.''} November 25.
\url{https://www.humanlayer.dev/blog/writing-a-good-claude-md}.

\bibitem[\citeproctext]{ref-liang2024repofuse}
{Liang, Ming, Xiaoheng Xie, Gehao Zhang, Xunjin Zheng, Peng Di, et al.}
2024. \emph{{REPOFUSE}: Repository-Level Code Completion with Fused Dual
Context}. \url{https://arxiv.org/abs/2402.14323}.

\bibitem[\citeproctext]{ref-aaif2025foundation}
Linux Foundation. 2025. {``{Linux Foundation} Announces the Formation of
the {Agentic AI Foundation}.''} December.
\url{https://www.linuxfoundation.org/press/linux-foundation-announces-the-formation-of-the-agentic-ai-foundation}.

\bibitem[\citeproctext]{ref-liu2024lost}
Liu, Nelson F., Kevin Lin, John Hewitt, et al. 2024. {``Lost in the
Middle: How Language Models Use Long Contexts.''} \emph{Transactions of
the Association for Computational Linguistics}.
\url{https://arxiv.org/abs/2307.03172}.

\bibitem[\citeproctext]{ref-liu2024graphcoder}
{Liu, Wei et al.} 2024. {``{GraphCoder}: Code Context Graph-Based
Retrieval for Repository-Level Completion.''} \emph{ASE 2024}.
\url{https://arxiv.org/abs/2406.07003}.

\bibitem[\citeproctext]{ref-phan2025repohyper}
{Phan, Huy N. et al.} 2025. {``{RepoHyper}: Search-Expand-Refine on
Semantic Graphs for Repository-Level Code Completion.''} \emph{FORGE
2025}. \url{https://arxiv.org/abs/2403.06095}.

\bibitem[\citeproctext]{ref-longcodebench2025}
Rando, Stefano, Luca Romani, Alessio Sampieri, et al. 2025.
\emph{{LongCodeBench}: Evaluating Coding LLMs at 1M Context Windows}.
arXiv preprint. \url{https://arxiv.org/abs/2505.07897}.

\bibitem[\citeproctext]{ref-shi2023distracted}
Shi, Freda, Xinyun Chen, Kanishka Misra, et al. 2023. {``Large Language
Models Can Be Easily Distracted by Irrelevant Context.''} \emph{ICML
2023}. \url{https://arxiv.org/abs/2302.00093}.

\bibitem[\citeproctext]{ref-shrivastava2023repofusion}
Shrivastava, Disha, Denis Kocetkov, Harm de Vries, Dzmitry Bahdanau, and
Torsten Scholak. 2023. \emph{{RepoFusion}: Training Code Models to
Understand Your Repository}. arXiv preprint.
\url{https://arxiv.org/abs/2306.10998}.

\bibitem[\citeproctext]{ref-sourcegraph2024scip}
Sourcegraph. 2024. {``{SCIP} --- {SourceCode Indexing Protocol}.''}
\url{https://github.com/sourcegraph/scip}.

\bibitem[\citeproctext]{ref-sourcegraph2025amp}
Sourcegraph. 2025. {``Why {Sourcegraph} and {Amp} Are Becoming
Independent Companies.''} December.
\url{https://sourcegraph.com/blog/why-sourcegraph-and-amp-are-becoming-independent-companies}.

\bibitem[\citeproctext]{ref-uniah2025}
{``{U-NIAH}: Unified {RAG} and {LLM} Evaluation for Long Context
Needle-in-a-Haystack.''} 2025. \emph{ACM Transactions on Information
Systems}. \url{https://arxiv.org/abs/2503.00353}.

\bibitem[\citeproctext]{ref-vercel2026nextevalsoss}
Vercel. 2026. {``{next-evals-oss}: Agent Evaluations for {Next.js}
Coding Tasks.''} \url{https://github.com/vercel/next-evals-oss}.

\bibitem[\citeproctext]{ref-latentspace2025claudecode}
Wang, Shawn (swyx), and Alessio Trotta. 2025. {``{Latent Space} Podcast:
{Claude Code} with Boris Cherny and Catherine Wu.''} May 7.
\url{https://www.latent.space/p/claude-code}.

\bibitem[\citeproctext]{ref-wang2025coderagbench}
{Wang, Zora Zhiruo et al.} 2025. {``{CodeRAG-Bench}: Can Retrieval
Augment Code Generation?''} \emph{NAACL Findings 2025}.
\url{https://arxiv.org/abs/2406.14497}.

\bibitem[\citeproctext]{ref-zed2025acp}
Zed Industries. 2025. {``Agent Client Protocol ({ACP}).''}
\url{https://zed.dev/acp}.

\bibitem[\citeproctext]{ref-zhang2023repocoder}
Zhang, Fengji, Bei Chen, Yue Zhang, et al. 2023. {``{RepoCoder}:
Repository-Level Code Completion Through Iterative Retrieval and
Generation.''} \emph{Proceedings of EMNLP 2023}.
\url{https://arxiv.org/abs/2303.12570}.

\bibitem[\citeproctext]{ref-zhang2025cast}
{Zhang, Yilin (Sherry) et al.} 2025. \emph{{cAST}: Enhancing Code RAG
with Structural Chunking via Abstract Syntax Tree}.
\url{https://arxiv.org/abs/2506.15655}.

\bibitem[\citeproctext]{ref-chen2024bgem3}
Chen, Jianlv, Shitao Xiao, Peitian Zhang, Kun Luo, Defu Lian, and Zheng
Liu. 2024. \emph{{BGE M3}-Embedding: Multi-Lingual,
Multi-Functionality, Multi-Granularity Text Embeddings Through
Self-Knowledge Distillation}. arXiv preprint.
\url{https://arxiv.org/abs/2402.03216}.

\bibitem[\citeproctext]{ref-formal2021splade}
Formal, Thibault, Benjamin Piwowarski, and Stéphane Clinchant. 2021.
{``{SPLADE}: Sparse Lexical and Expansion Model for First Stage
Ranking.''} \emph{SIGIR 2021}. \url{https://arxiv.org/abs/2107.05720}.

\bibitem[\citeproctext]{ref-fu2025absencebench}
Fu, Harvey Yiyun, Aryan Shrivastava, Jared Moore, Peter West, Chenhao
Tan, and Ari Holtzman. 2025. {``AbsenceBench: Language Models Can't Tell
What's Missing.''} \emph{NeurIPS 2025 Datasets and Benchmarks Track}.
\url{https://arxiv.org/abs/2506.11440}.

\bibitem[\citeproctext]{ref-hong2025contextrot}
Hong, Kelly, Anton Troynikov, and Jeff Huber. 2025. {``Context Rot: How
Increasing Input Tokens Impacts LLM Performance.''} Chroma Technical
Report. \url{https://trychroma.com/research/context-rot}.

\bibitem[\citeproctext]{ref-husain2019codesearchnet}
Husain, Hamel, Ho-Hsiang Wu, Tiferet Gazit, Miltiadis Allamanis, and
Marc Brockschmidt. 2019. \emph{CodeSearchNet Challenge: Evaluating the
State of Semantic Code Search}. arXiv preprint.
\url{https://arxiv.org/abs/1909.09436}.

\bibitem[\citeproctext]{ref-jiang2025cosil}
Jiang, Zhonghao, Xiaoxue Ren, Meng Yan, Wei Jiang, Yong Li, and Zhongxin
Liu. 2025. \emph{{CoSIL}: Software Issue Localization via LLM-Driven
Code Repository Graph Searching}. arXiv preprint.
\url{https://arxiv.org/abs/2503.22424}.

\bibitem[\citeproctext]{ref-khattab2020colbert}
Khattab, Omar, and Matei Zaharia. 2020. {``{ColBERT}: Efficient and
Effective Passage Search via Contextualized Late Interaction over
BERT.''} \emph{Proceedings of SIGIR 2020}, 39--48.
\url{https://doi.org/10.1145/3397271.3401075}.

\bibitem[\citeproctext]{ref-lee2024mxbaiembed}
Lee, Sean, Aamir Shakir, Darius Koenig, and Julius Lipp. 2024a. {``Open
Source Strikes Bread -- New Fluffy Embedding Model.''} Mixedbread.
\url{https://www.mixedbread.com/blog/mxbai-embed-large-v1}.

\bibitem[\citeproctext]{ref-lee2024mxbaicolbert}
Lee, Sean, Aamir Shakir, Darius Koenig, and Julius Lipp. 2024b.
{``ColBERTus Maximus -- Introducing mxbai-colbert-large-v1.''}
Mixedbread. \url{https://www.mixedbread.com/blog/mxbai-colbert-large-v1}.

\bibitem[\citeproctext]{ref-li2024coir}
Li, Xiangyang, Kuicai Dong, Yi Quan Lee, Wei Xia, Yichun Yin, Hao Zhang,
Yong Liu, Yasheng Wang, and Ruiming Tang. 2024. \emph{{CoIR}: A
Comprehensive Benchmark for Code Information Retrieval Models}. arXiv
preprint. \url{https://arxiv.org/abs/2407.02883}.

\bibitem[\citeproctext]{ref-mixedbread2026mgrep}
Mixedbread AI. 2026. {``mgrep: Semantic Search for Agents.''}
\url{https://www.mgrep.dev/}.

\bibitem[\citeproctext]{ref-modarressi2025nolima}
Modarressi, Ali, Hanieh Deilamsalehy, Franck Dernoncourt, Trung Bui,
Ryan A. Rossi, Seunghyun Yoon, and Hinrich Schütze. 2025.
{``NoLiMa: Long-Context Evaluation Beyond Literal Matching.''}
\emph{International Conference on Machine Learning}.
\url{https://arxiv.org/abs/2502.05167}.

\bibitem[\citeproctext]{ref-mohsenimofidi2025contextoss}
Mohsenimofidi, Seyedmoein, Matthias Galster, Christoph Treude, and
Sebastian Baltes. 2025. \emph{Context Engineering for AI Agents in
Open-Source Software}. arXiv preprint.
\url{https://arxiv.org/abs/2510.21413}.

\bibitem[\citeproctext]{ref-muennighoff2023mteb}
Muennighoff, Niklas, Nouamane Tazi, Loic Magne, and Nils Reimers. 2023.
{``{MTEB}: Massive Text Embedding Benchmark.''} \emph{Proceedings of
EACL 2023}. \url{https://arxiv.org/abs/2210.07316}.

\bibitem[\citeproctext]{ref-robertson2009bm25}
Robertson, Stephen, and Hugo Zaragoza. 2009. {``The Probabilistic
Relevance Framework: {BM25} and Beyond.''} \emph{Foundations and Trends
in Information Retrieval} 3 (4): 333--389.
\url{https://doi.org/10.1561/1500000019}.

\bibitem[\citeproctext]{ref-santhanam2022colbertv2}
Santhanam, Keshav, Omar Khattab, Jon Saad-Falcon, Christopher Potts, and
Matei Zaharia. 2022a. {``{ColBERTv2}: Effective and Efficient Retrieval
via Lightweight Late Interaction.''} \emph{NAACL 2022}.
\url{https://aclanthology.org/2022.naacl-main.272/}.

\bibitem[\citeproctext]{ref-santhanam2022plaid}
Santhanam, Keshav, Omar Khattab, Christopher Potts, and Matei Zaharia.
2022b. \emph{{PLAID}: An Efficient Engine for Late Interaction
Retrieval}. arXiv preprint. \url{https://arxiv.org/abs/2205.09707}.

\bibitem[\citeproctext]{ref-thakur2021beir}
Thakur, Nandan, Nils Reimers, Andreas Rücklé, Abhishek Srivastava, and
Iryna Gurevych. 2021. {``{BEIR}: A Heterogeneous Benchmark for
Zero-Shot Evaluation of Information Retrieval Models.''} \emph{NeurIPS
2021 Datasets and Benchmarks Track}.
\url{https://openreview.net/forum?id=wCu6T5xFjeJ}.

\bibitem[\citeproctext]{ref-thakur2025freshstack}
Thakur, Nandan, Jimmy Lin, Sam Havens, Michael Carbin, Omar Khattab, and
Andrew Drozdov. 2025. {``FreshStack: Building Realistic Benchmarks for
Evaluating Retrieval on Technical Documents.''} \emph{NeurIPS 2025
Datasets and Benchmarks Track}. \url{https://arxiv.org/abs/2504.13128}.

\end{CSLReferences}




\end{document}
